% Chapter 3, Section 3 _Linear Algebra_ Jim Hefferon
%  http://joshua.smcvt.edu/linearalgebra
%  2001-Jun-12
\section{Menghitung Peta Linear}
Bagian sebelumnya menunjukkan bahwa
suatu peta linear ditentukan oleh aksinya pada suatu basis.
Persamaan
\begin{equation*}
  h(\vec{v})
  =h(c_1\cdot\vec{\beta}_1+\dots+c_n\cdot\vec{\beta}_n)
  =c_1\cdot h(\vec{\beta}_1)+\dots +c_n\cdot h(\vec{\beta}_n)
\tag*{}\end{equation*}
menjelaskan cara memperoleh nilai peta pada sebarang vektor $\vec{v}$
dengan bermula 
dari nilai peta pada vektor-vektor $\vec{\beta}_i$ dalam suatu basis 
dan memperluasnya secara linear.
% We just need to  
% find the $c$'s to express $\vec{v}$ with respect to the basis.

Bagian ini memberikan suatu skema praktis berbasis matriks
untuk menggunakan representasi 
\( h(\vec{\beta}_1) \), \ldots, \( h(\vec{\beta}_n) \)
guna menghitung, dari representasi suatu vektor dalam domain 
$\rep{\vec{v}}{B}$,
representasi bayangan vektor tersebut dalam kodomain 
$\rep{h(\vec{v})}{D}$.














\subsection{Merepresentasikan Peta Linear dengan Matriks}
\index{homomorfisme!matriks yang merepresentasikan|(}
\begin{example}  \label{ex:TypLinMapRepByMat}
Untuk ruang $\Re^2$ dan $\Re^3$, tetapkan basis-basis berikut.
\begin{equation*}
   B=\sequence{
               \colvec[r]{2 \\ 0},
               \colvec[r]{1 \\ 4}}
   \qquad
   D=\sequence{
               \colvec[r]{1 \\ 0 \\ 0},
               \colvec[r]{0 \\ -2 \\ 0},
               \colvec[r]{1 \\ 0 \\ 1}}
\end{equation*}
Tinjau peta $\map{h}{\Re^2}{\Re^3}$ yang ditentukan oleh pengaitan berikut.
\begin{equation*}
  \colvec[r]{2 \\ 0}
    \mapsunder{h}
  \colvec[r]{1 \\ 1 \\ 1}
  \qquad
  \colvec[r]{1 \\ 4}
    \mapsunder{h}
  \colvec[r]{1 \\ 2 \\ 0}
\end{equation*}
Untuk menghitung aksi peta ini pada sebarang vektor dalam domain,
mula-mula kita merepresentasikan vektor $h(\vec{\beta}_1)$
\begin{equation*}
  \colvec[r]{1 \\ 1 \\ 1}=
         0\colvec[r]{1 \\ 0 \\ 0}
         -\frac{1}{2}\colvec[r]{0 \\ -2 \\ 0}
         +1\colvec[r]{1 \\ 0 \\ 1}
  \qquad
   \rep{ h(\vec{\beta}_1) }{D}=\colvec[r]{0 \\ -1/2 \\ 1}_D           
\end{equation*}
dan $h(\vec{\beta}_2)$.
\begin{equation*}
  \colvec[r]{1 \\ 2 \\ 0}=
        1\colvec[r]{1 \\ 0 \\ 0}
        -1\colvec[r]{0 \\ -2 \\ 0}
        +0\colvec[r]{1 \\ 0 \\ 1}
  \qquad
  \rep{ h(\vec{\beta}_2) }{D}=\colvec[r]{1 \\ -1 \\ 0}_D
\end{equation*}
Dengan ini, untuk sebarang anggota $\vec{v}$ dari domain,
kita dapat menghitung $h(\vec{v})$.
\begin{align*}
  h(\vec{v})
  &=h(c_1\cdot \colvec[r]{2 \\ 0}+c_2\cdot \colvec[r]{1 \\ 4})   \\
  &=c_1\cdot h(\colvec[r]{2 \\ 0})+c_2\cdot h(\colvec[r]{1 \\ 4}) \\
  &=c_1\cdot (
        0\colvec[r]{1 \\ 0 \\ 0}
        \!-\frac{1}{2}\colvec[r]{0 \\ -2 \\ 0}
        \!+1\colvec[r]{1 \\ 0 \\ 1}\!  )
  +c_2\cdot (
        1\colvec[r]{1 \\ 0 \\ 0}
        \!-1\colvec[r]{0 \\ -2 \\ 0}
        \!+0\colvec[r]{1 \\ 0 \\ 1}\!  )      \\
  &=(0c_1+1c_2)\cdot \colvec[r]{1 \\ 0 \\ 0}
   +(-\frac{1}{2}c_1-1c_2)\cdot \colvec[r]{0 \\ -2 \\ 0}
   +(1c_1+0c_2)\cdot \colvec[r]{1 \\ 0 \\ 1}
\end{align*}
Jadi,
\begin{center}
  jika $\rep{\vec{v}}{B}=\colvec{c_1 \\ c_2}$
  maka $\rep{\,h(\vec{v})\,}{D}
  =\colvec{0c_1+1c_2 \\ -(1/2)c_1-1c_2 \\ 1c_1+0c_2}$.
\end{center}
Sebagai contoh, 
\begin{center}
  karena $\rep{\colvec[r]{4 \\ 8}}{B}=\colvec[r]{1 \\ 2}_B$,
  kita memperoleh
  $\rep{\,h(\colvec[r]{4 \\ 8})\,}{D}
   =\colvec[r]{2 \\ -5/2 \\ 1}$.
\end{center}
\end{example}

Kita menyatakan perhitungan seperti di atas dengan notasi matriks.
\begin{equation*}
    \begin{mat}[r]
      0             &1  \\
      -1/2          &-1  \\
      1             &0
    \end{mat}_{B,D}
  \colvec{c_1 \\ c_2}_B
  =
  \colvec{0c_1+1c_2 \\ (-1/2)c_1-1c_2 \\ 1c_1+0c_2}_D
\end{equation*}
Di tengah terdapat argumen $\vec{v}$ bagi peta tersebut, 
yang direpresentasikan terhadap basis domain $B$
oleh vektor kolom dengan komponen $c_1$ dan $c_2$.
Di sebelah kanan terdapat nilai peta pada argumen itu, $h(\vec{v})$,
yang direpresentasikan terhadap basis kodomain $D$.
Matriks di sebelah kiri adalah unsur yang baru.
Kita akan menggunakannya untuk merepresentasikan peta dan memandang 
persamaan di atas sebagai representasi penerapan peta pada vektor tersebut.

Matriks itu terdiri atas koefisien-koefisien dari vektor di sebelah kanan,
$0$ dan $1$ dari baris pertama, $-1/2$ dan $-1$ dari
baris kedua, serta $1$ dan $0$ dari baris ketiga.
Dengan kata lain, kita menyusunnya dengan menyejajarkan vektor-vektor 
yang merepresentasikan $h(\vec{\beta}_i)$. 
\begin{equation*}
  \left(\begin{array}{c|c}
     \vdots                         &\vdots    \\
     \rep{\,h(\vec{\beta}_1)\,}{D}  &\rep{\,h(\vec{\beta}_2)\,}{D}  \\
     \vdots                         &\vdots
  \end{array}\right)
\end{equation*}

\begin{definition} \label{def:MatRepMap}
%<*df:MatRepMap>
Andaikan \( V \) dan \( W \) adalah ruang vektor berdimensi \( n \) dan
\( m \) dengan basis \( B \) dan \( D \),
dan \( \map{h}{V}{W} \) adalah peta linear.
Jika
\begin{equation*}
  \rep{h( \vec{\beta}_1 )}{D}
  =
  \colvec{h_{1,1} \\ h_{2,1} \\ \vdots \\ h_{m,1}}_D
  \quad\ldots\quad
  \rep{h( \vec{\beta}_n )}{D}
  =
  \colvec{h_{1,n} \\ h_{2,n} \\ \vdots \\ h_{m,n}}_D
\end{equation*}
maka 
\begin{equation*}
  \rep{h}{B,D}=\generalmatrix{h}{n}{m}_{B,D}
\end{equation*}
adalah \definend{representasi matriks dari \( h \) terhadap \( B, D \)}.%
\index{representasi!matriks}\index{matriks!representasi}%
\index{homomorfisme!matriks yang merepresentasikan}
%</df:MatRepMap>
\end{definition}

\noindent Dalam matriks tersebut, banyaknya kolom~$n$ adalah 
dimensi domain peta, sedangkan  
banyaknya baris~$m$ adalah dimensi kodomain.

\begin{remark}
Seperti notasi untuk representasi suatu vektor, notasi $\mbox{Rep}_{B,D}$
di sini bukan notasi baku.
Alternatif yang paling umum adalah $[h]_{B,D}$. 
\end{remark}

Kita menggunakan huruf kecil untuk peta, 
huruf besar untuk matriks,
dan huruf kecil lagi untuk entri-entri matriks.
Jadi, untuk peta \( h \), matriks yang merepresentasikannya adalah \( H \), dengan
entri \( h_{i,j} \).

\begin{example}  \label{ex:PolyOneToRThree}
Jika \( \map{h}{\Re^3}{\polyspace_1} \) adalah
\begin{equation*}
  \colvec{a_1 \\ a_2 \\ a_3}
     \mapsunder{h}
   (2a_1+a_2)+(-a_3)x
\end{equation*}
maka, dengan
\begin{equation*}
  B=
  \sequence{\colvec[r]{0 \\ 0 \\ 1},
            \colvec[r]{0 \\ 2 \\ 0},
            \colvec[r]{2 \\ 0 \\ 0} }
  \qquad
  D=
  \sequence{1+x,-1+x}
\end{equation*}
aksi \( h \) pada \( B \) adalah sebagai berikut.
\begin{equation*}
  \colvec[r]{0 \\ 0 \\ 1}\mapsunder{h}-x
  \qquad \colvec[r]{0 \\ 2 \\ 0}\mapsunder{h}2
  \qquad \colvec[r]{2 \\ 0 \\ 0}\mapsunder{h}4
\end{equation*}
Perhitungan sederhana
\begin{equation*}
  \rep{-x}{D}=\colvec[r]{-1/2 \\ -1/2}_D
  \quad \rep{2}{D}=\colvec[r]{1 \\ -1}_D 
  \quad \rep{4}{D}=\colvec[r]{2 \\ -2}_D 
\end{equation*}
menunjukkan bahwa inilah matriks yang merepresentasikan $h$ terhadap basis-basis tersebut.
\begin{equation*}
  \rep{h}{B,D}
   =
   \begin{mat}[r]
     -1/2  &1   &2  \\
     -1/2  &-1  &-2
   \end{mat}_{B,D}
\end{equation*}
\end{example}

\begin{theorem} \label{th:MatMultRepsFuncAppl}
%<*th:MatMultRepsFuncAppl>
Andaikan \( V \) dan \( W \) adalah ruang vektor
berdimensi \( n \) dan \( m \)
dengan basis \( B \) dan \( D \),
dan \( \map{h}{V}{W} \) adalah peta linear.
Jika \( h \) direpresentasikan oleh
\begin{equation*}
  \rep{h}{B,D}=\generalmatrix{h}{n}{m}_{B,D}
\end{equation*}
dan \( \vec{v}\in V \) direpresentasikan oleh
\begin{equation*}
  \rep{\vec{v}}{B}=\colvec{c_1 \\ c_2 \\ \vdots \\ c_n}_B
\end{equation*}
maka representasi bayangan $\vec{v}$ adalah sebagai berikut.
\begin{equation*}
  \rep{\, h(\vec{v}) \,}{D}
  =
  \colvec{h_{1,1}c_1+h_{1,2}c_2+\dots+h_{1,n}c_n \\
          h_{2,1}c_1+h_{2,2}c_2+\dots+h_{2,n}c_n \\
          \vdots \\
          h_{m,1}c_1+h_{m,2}c_2+\dots+h_{m,n}c_n}_D
\end{equation*}
%</th:MatMultRepsFuncAppl>
\end{theorem}

\begin{proof}
%<*pf:MatMultRepsFuncAppl>
Ini memformalkan \nearbyexample{ex:TypLinMapRepByMat}.
Lihat \nearbyexercise{exer:MatVecMultRepLinMap}.
%</pf:MatMultRepsFuncAppl>
\end{proof}

\begin{definition} \label{def:MatrixVecProd}
%<*df:MatrixVecProd>
\definend{Hasil kali matriks-vektor}\index{perkalian!matriks-vektor}%
\index{matriks!hasil kali matriks-vektor}
dari suatu matriks \( \nbym{m}{n} \) dan suatu
vektor \( \nbym{n}{1} \) adalah sebagai berikut.
\begin{equation*}
  \generalmatrix{a}{n}{m}
  \colvec{c_1 \\ \vdots \\ c_n}
  =
  \colvec{a_{1,1}c_1+\dots+a_{1,n}c_n \\
             a_{2,1}c_1+\dots+a_{2,n}c_n \\
             \vdots \\ 
             a_{m,1}c_1+\dots+a_{m,n}c_n}
\end{equation*}
%</df:MatrixVecProd>
\end{definition}

% the product of the matrix $\rep{h}{B,D}$ with the vector $\rep{\vec{v}}{B}$ 
% is the vector $\rep{h(\vec{v})}{D}$.  
Singkatnya, 
penerapan suatu peta linear direpresentasikan oleh hasil kali matriks-vektor 
antara representasi peta dan representasi vektor.

\begin{remark}  \label{rem:NotSurprising}
\nearbytheorem{th:MatMultRepsFuncAppl} tidaklah mengejutkan,
karena kita memilih representasi matriks dalam \nearbydefinition{def:MatRepMap}
tepat agar teorema tersebut benar\Dash
jika teorema itu tidak benar, kita akan menyesuaikan definisinya
agar menjadi benar.
Meskipun demikian, kita tetap memerlukan verifikasi. 
\end{remark}

\begin{example}
Untuk matriks dari \nearbyexample{ex:PolyOneToRThree},
kita dapat menghitung ke mana peta tersebut mengirim vektor berikut.
\begin{equation*}
  \vec{v}=\colvec[r]{4 \\ 1 \\ 0}
\end{equation*}
Terhadap basis domain $B$, representasi vektor ini adalah 
\begin{equation*}
  \rep{\vec{v}}{B}=\colvec[r]{0 \\ 1/2 \\ 2}_B
\end{equation*}
sehingga hasil kali matriks-vektor memberikan 
representasi nilai $h(\vec{v})$ terhadap
basis kodomain $D$.
\begin{align*}
  \rep{h(\vec{v})}{D}
  &=\begin{mat}[r]
      -1/2  &1   &2  \\
      -1/2  &-1  &-2
    \end{mat}_{B,D}
    \colvec[r]{0 \\ 1/2 \\ 2}_B                            \\
  &=\colvec{(-1/2)\cdot 0+1\cdot (1/2) + 2\cdot 2 \\ 
          (-1/2)\cdot 0-1\cdot (1/2) - 2\cdot 2}_D
  =\colvec[r]{9/2 \\ -9/2}_D
\end{align*}
Untuk mencari $h(\vec{v})$ itu sendiri, bukan representasinya,
hitung $(9/2)(1+x)-(9/2)(-1+x)=9$.
\end{example}

\begin{example}
Misalkan \( \map{\pi}{\Re^3}{\Re^2} \) adalah proyeksi pada bidang \( xy \).
Untuk memberikan matriks yang merepresentasikan peta ini, mula-mula kita tetapkan beberapa basis. 
\begin{equation*}
  B=\sequence{
              \colvec[r]{1 \\ 0 \\ 0},
              \colvec[r]{1 \\ 1 \\ 0},
              \colvec[r]{-1 \\ 0 \\ 1} }
  \qquad
  D=\sequence{
              \colvec[r]{2 \\ 1},
              \colvec[r]{1 \\ 1} }
\end{equation*}
Untuk setiap vektor dalam basis domain, cari bayangannya di bawah peta tersebut. 
\begin{equation*}
  \colvec[r]{1 \\ 0 \\ 0}\mapsunder{\pi}\colvec[r]{1 \\ 0}
  \quad
  \colvec[r]{1 \\ 1 \\ 0}\mapsunder{\pi}\colvec[r]{1 \\ 1}
  \quad
  \colvec[r]{-1 \\ 0 \\ 1}\mapsunder{\pi}\colvec[r]{-1 \\ 0}
\end{equation*}
Kemudian cari representasi setiap bayangan terhadap basis
kodomain. 
\begin{equation*}
  \rep{\colvec[r]{1 \\ 0}}{D}=\colvec[r]{1 \\ -1}
  \quad
  \rep{\colvec[r]{1 \\ 1}}{D}=\colvec[r]{0 \\ 1}
  \quad
  \rep{\colvec[r]{-1 \\ 0}}{D}=\colvec[r]{-1 \\ 1}
\end{equation*}
Terakhir, menyejajarkan representasi-representasi ini memberikan matriks yang merepresentasikan 
\( \pi \) terhadap \( B,D \).
\begin{equation*}
    \rep{\pi}{B,D}
    =\begin{mat}[r]
      1  &0  &-1  \\
      -1 &1  &1
    \end{mat}_{B,D}
\end{equation*}
Kita dapat mengilustrasikan \nearbytheorem{th:MatMultRepsFuncAppl} dengan menghitung
hasil kali matriks-vektor yang merepresentasikan aksi 
peta proyeksi berikut.
\begin{equation*}
  \pi(
     \colvec[r]{2 \\ 2 \\ 1}
     )=\colvec[r]{2 \\ 2}
\end{equation*}
Representasikan vektor domain 
terhadap basis domain
\begin{equation*}
   \rep{\colvec[r]{2 \\ 2 \\ 1}}{B}=
        \colvec[r]{1 \\ 2 \\ 1}_B
\end{equation*}
untuk memperoleh hasil kali matriks-vektor berikut.
\begin{equation*}
   \rep{ \,\pi(\colvec[r]{2 \\ 2 \\ 1})\,}{D}=
      \begin{mat}[r]
        1  &0  &-1  \\
        -1 &1  &1
      \end{mat}_{B,D}
   \colvec[r]{1 \\ 2 \\ 1}_B
   =
   \colvec[r]{0 \\ 2}_D
\end{equation*}
Menguraikannya menjadi kombinasi linear vektor-vektor dari
\( D \) 
\begin{equation*}
   0\cdot\colvec[r]{2 \\ 1}
   +2\cdot\colvec[r]{1 \\ 1}
   =
   \colvec[r]{2 \\ 2}
\end{equation*}
memastikan bahwa aksi peta memang
tercermin dalam operasi matriks tersebut.
Kadang-kadang kita akan memadatkan ketiga persamaan tersaji ini menjadi satu. 
\begin{equation*}
  \colvec[r]{2 \\ 2 \\ 1}=\colvec[r]{1 \\ 2 \\ 1}_B
    \;\overset{h}{\underset{H}{\longmapsto}}\;
  \colvec[r]{0 \\ 2}_D=\colvec[r]{2 \\ 2}    
\end{equation*}
\end{example}

Kini kita memiliki dua cara untuk menghitung efek proyeksi,
yakni rumus langsung yang membuang komponen ketiga setiap vektor bertinggi tiga
untuk membentuk vektor bertinggi dua, 
dan rumus di atas yang menggunakan representasi serta perkalian 
matriks-vektor.
Cara kedua mungkin tampak rumit dibandingkan cara pertama,
tetapi memiliki kelebihan.
Contoh berikut menunjukkan bahwa untuk beberapa peta, skema baru ini menyederhanakan
rumusnya.

\begin{example} \label{exam:RepsOfRigidPlaneMaps}
Untuk merepresentasikan peta rotasi\index{rotasi (atau perputaran)!representasi}
$\map{t_{\theta}}{\Re^2}{\Re^2}$ yang 
memutar semua vektor di bidang berlawanan arah jarum jam sebesar sudut $\theta$,
\begin{center}  \small
  \includegraphics{map/mp/ch3.15}
\end{center}
kita mulai dengan menetapkan basis standar
$\stdbasis_2$ sebagai basis domain maupun kodomain.
Kemudian cari bayangan setiap vektor
dalam basis domain di bawah peta tersebut.
\begin{equation*}
  \colvec[r]{1 \\ 0}\mapsunder{t_\theta}\colvec{\cos\theta \\ \sin\theta}
  \qquad
  \colvec[r]{0 \\ 1}\mapsunder{t_\theta}\colvec{-\sin\theta \\ \cos\theta}
  \tag{$*$}
\end{equation*}
Representasikan bayangan-bayangan ini terhadap basis kodomain.
Karena basis ini adalah $\stdbasis_2$, vektor merepresentasikan dirinya sendiri.
Sejajarkan representasi-representasi tersebut untuk memperoleh matriks yang merepresentasikan peta.
\begin{equation*}
  \rep{t_\theta}{\stdbasis_2,\stdbasis_2}
  =
  \begin{mat}
    \cos\theta  &-\sin\theta \\
    \sin\theta  &\cos\theta
  \end{mat}
\end{equation*}
Kelebihan skema ini adalah bahwa kita memperoleh rumus bagi 
bayangan sebarang vektor 
hanya dengan mengetahui, dalam~($*$), cara
merepresentasikan bayangan kedua vektor basis. 
Sebagai contoh, di sini kita merotasi suatu vektor dengan $\theta=\pi/6$.
\begin{equation*}
  \colvec[r]{3  \\ -2} = \colvec[r]{3  \\ -2}_{\stdbasis_2}\!\mapsunder{t_{\pi/6}}\;
  \begin{mat}[r]
    \sqrt{3}/2  &-1/2  \\
     1/2        &\sqrt{3}/2
  \end{mat}
  \colvec[r]{3  \\ -2}
  \approx
  \colvec[r]{3.598 \\ -0.232}_{\stdbasis_2}
  \!\!=  
  \colvec[r]{3.598 \\ -0.232}
\end{equation*}
Secara lebih umum, kita memiliki rumus untuk rotasi sebesar $\theta=\pi/6$.
\begin{equation*}
  \colvec[r]{x  \\ y} \mapsunder{t_{\pi/6}}\;
  \begin{mat}[r]
    \sqrt{3}/2  &-1/2  \\
     1/2        &\sqrt{3}/2
  \end{mat}
  \colvec[r]{x  \\ y}
  =
  \colvec[r]{(\sqrt{3}/2)x-(1/2)y \\ (1/2)x+(\sqrt{3}/2)y }
\end{equation*}
\end{example}

\begin{example}
Dalam definisi hasil kali matriks-vektor,
lebar matriks sama dengan tinggi vektor.
Karena itu, hasil kali berikut tidak terdefinisi. 
\begin{equation*}
    \begin{mat}[r]
      1  &0  &0  \\
      4  &3  &1
    \end{mat}
  \colvec[r]{1 \\ 0}
\end{equation*}
Ada alasan mengapa hasil kali ini tidak terdefinisi:~matriks berlebar tiga
merepresentasikan peta dengan domain berdimensi tiga, sedangkan 
vektor bertinggi dua merepresentasikan anggota ruang berdimensi dua.
Jadi, vektor itu tidak mungkin berada dalam domain peta.
\end{example}

Tidak ada dalam \nearbydefinition{def:MatrixVecProd} yang mengharuskan kita memandang
hasil kali matriks-vektor melalui representasi.
Kita dapat memperoleh beberapa wawasan dengan berfokus pada cara entri-entri berpadu.

Cara yang baik untuk memandang hasil kali matriks-vektor adalah bahwa hasil itu dibentuk dari 
hasil kali titik antara baris-baris matriks dan vektor kolom.
\begin{equation*}
    \begin{mat}
               &\vdots                         \\
      a_{i,1}  &a_{i,2}  &\ldots   &a_{i,n}    \\
               &\vdots
    \end{mat}
  \colvec{c_1 \\ c_2 \\ \vdots \\ c_n}
  =
  \colvec{\vdots \\ a_{i,1}c_1+a_{i,2}c_2+\cdots+a_{i,n}c_n \\ \vdots}
\end{equation*}
Jika dipandang baris demi baris seperti ini,
operasi baru ini memperumum hasil kali titik.

Kita juga dapat memandang operasi ini kolom demi kolom.
\begin{align*}
           \generalmatrix{h}{n}{m}
           \colvec{c_1 \\ c_2 \\ \vdots \\ c_n}
  &=\colvec{h_{1,1}c_1+h_{1,2}c_2+\dots+h_{1,n}c_n \\
               h_{2,1}c_1+h_{2,2}c_2+\dots+h_{2,n}c_n \\
               \vdots \\
               h_{m,1}c_1+h_{m,2}c_2+\dots+h_{m,n}c_n}    \\
  &=c_1\colvec{h_{1,1} \\ h_{2,1} \\ \vdots \\ h_{m,1}}
%   +c_2\colvec{h_{1,2} \\ h_{2,2} \\ \vdots \\ h_{m,2}}
   +\dots
   +c_n\colvec{h_{1,n} \\ h_{2,n} \\ \vdots \\ h_{m,n}}
\end{align*}
Hasilnya adalah
kolom-kolom matriks yang diboboti oleh entri-entri vektor.

\begin{example}
\begin{equation*}
    \begin{mat}[r]
      1  &0  &-1  \\
      2  &0  &3
    \end{mat}
  \colvec[r]{2 \\ -1 \\ 1}
  =
  2\colvec[r]{1 \\ 2}
  -1\colvec[r]{0 \\ 0}
  +1\colvec[r]{-1 \\ 3}
  =
  \colvec[r]{1 \\ 7}
\end{equation*}
\end{example}

Cara memandang hasil kali matriks-vektor ini
membawa kita kembali pada tujuan yang dinyatakan di awal bagian ini, yakni menghitung
\( h(c_1\vec{\beta}_1+\dots+c_n\vec{\beta}_n) \)
sebagai
\( c_1h(\vec{\beta}_1)+\dots+c_nh(\vec{\beta}_n) \).

Kita memulai bagian ini
dengan mencatat bahwa kesamaan keduanya memungkinkan kita menghitung aksi 
$h$ pada sebarang
argumen hanya dengan mengetahui $h(\vec{\beta}_1)$, \ldots, $h(\vec{\beta}_n)$.
Kita telah mengembangkannya menjadi suatu skema untuk
menghitung aksi peta dengan mengambil 
hasil kali matriks-vektor antara matriks yang merepresentasikan 
peta dan vektor yang merepresentasikan argumen.
Dengan cara ini, terhadap sebarang basis, setiap peta linear memiliki 
representasi matriks.
Subbagian berikut akan menunjukkan kebalikannya: jika kita menetapkan basis, maka 
setiap matriks memiliki peta linear yang terkait.


\begin{exercises}
  \recommended \item  
    Kalikan matriks
    \begin{equation*}
      \begin{mat}[r]
        1  &3  &1  \\
        0  &-1 &2  \\
        1  &1  &0
      \end{mat}
    \end{equation*}
    dengan setiap vektor, atau nyatakan ``tidak terdefinisi.''
    \begin{exparts*}
      \partsitem \( \colvec[r]{2 \\ 1 \\ 0} \)
      \partsitem \( \colvec[r]{-2 \\ -2} \)
      \partsitem \( \colvec[r]{0 \\ 0 \\ 0} \)
    \end{exparts*}
    \begin{answer}
      \begin{exparts*}
        \partsitem \(
           \colvec{1\cdot 2+3\cdot 1+1\cdot 0         \\
                        0\cdot 2+(-1)\cdot 1+2\cdot 0 \\
                        1\cdot 2+1\cdot 1+0\cdot 0}
           =\colvec[r]{5 \\ -1 \\ 3}   \)
        \partsitem Tidak terdefinisi.
        \partsitem \(  \colvec[r]{0 \\ 0 \\ 0}  \)
      \end{exparts*}  
    \end{answer}
  \item Kalikan matriks ini dengan setiap vektor atau nyatakan ``tidak terdefinisi.''
    \begin{equation*}
      \begin{mat}
        3  &1  \\
        2  &4
      \end{mat}
    \end{equation*}
    \begin{exparts*}
      \partsitem $\colvec{1 \\ 2  \\ 1}$
      \partsitem $\colvec{0 \\ -1}$
      \partsitem $\colvec{0 \\ 0 \\ 0}$
    \end{exparts*}
    \begin{answer}
      \begin{exparts}
        \partsitem Ini tidak terdefinisi.
        \partsitem Ini terdefinisi.
          \begin{equation*}
            \begin{mat}
              3  &1  \\
              2  &4
            \end{mat}
            \colvec{0 \\ -1}
            =
            \colvec{-1 \\ -4}
          \end{equation*}
       \partsitem Tidak terdefinisi.
      \end{exparts}
    \end{answer}
  \item 
    Lakukan setiap perkalian matriks-vektor berikut, jika mungkin.
    \begin{exparts*}
      \partsitem $\begin{mat}[r]
                    2  &1  \\
                    3  &-1/2
                  \end{mat}
                  \colvec[r]{4  \\ 2}$
      \partsitem $\begin{mat}[r]
                    1  &1  &0 \\
                    -2 &1  &0
                  \end{mat}
                  \colvec[r]{1 \\ 3 \\ 1}$
      \partsitem $\begin{mat}[r]
                    1  &1  \\
                    -2  &1
                  \end{mat}
                  \colvec[r]{1 \\ 3 \\ 1}$
    \end{exparts*}
    \begin{answer}
      \begin{exparts*}
        \partsitem $\colvec{2\cdot 4 +1\cdot 2 \\
                            3\cdot 4-(1/2)\cdot 2}
                   =\colvec[r]{10 \\ 11}$
        \partsitem $\colvec[r]{4 \\ 1}$
        \partsitem Tidak terdefinisi.
      \end{exparts*}
    \end{answer}
  \item  
    Persamaan matriks ini menyatakan suatu sistem linear.
    Selesaikan.
    \begin{equation*}
      \begin{mat}[r]
        2  &1  &1  \\
        0  &1  &3  \\
        1  &-1 &2
      \end{mat}
      \colvec{x \\ y \\ z}
      =\colvec[r]{8 \\ 4 \\ 4}
    \end{equation*}
    \begin{answer}
      Perkalian matriks-vektor menghasilkan suatu sistem linear.
      \begin{equation*}
        \begin{linsys}{3}
          2x  &+  &y  &+  &z  &=  &8  \\
              &   &y  &+  &3z &=  &4  \\
           x  &-  &y  &+  &2z &=  &4 
        \end{linsys}
      \end{equation*}
      Reduksi Gauss menunjukkan bahwa \( z=1 \), \( y=1 \), dan \( x=3 \).  
    \end{answer}
  \recommended \item 
    Untuk suatu homomorfisme dari \( \polyspace_2 \) ke \( \polyspace_3 \) yang memetakan
    \begin{equation*}
      1\mapsto 1+x,
      \quad
      x\mapsto 1+2x,
      \quad\text{dan}\quad
      x^2\mapsto x-x^3
    \end{equation*}
    ke manakah \( 1-3x+2x^2 \) dipetakan?
    \begin{answer}
      Berikut dua cara memperoleh jawabannya.

      Pertama, jelas bahwa $1-3x+2x^2=1\cdot 1-3\cdot x+2\cdot x^2$, sehingga kita dapat
      menerapkan sifat umum pelestarian kombinasi untuk memperoleh
      $h(1-3x+2x^2)
       =h(1\cdot 1-3\cdot x+2\cdot x^2) 
       =1\cdot h(1)-3\cdot h(x)+2\cdot h(x^2) 
       =1\cdot (1+x)-3\cdot (1+2x)+2\cdot (x-x^3) 
       =-2-3x-2x^3$.

      Cara lainnya menggunakan skema perhitungan yang dikembangkan dalam subbagian ini.
      Karena kita mengetahui ke mana unsur-unsur ruang ini dipetakan, kita meninjau
      basis \( B=\sequence{1,x,x^2} \) untuk domain.
      Secara sebarang, kita dapat mengambil \( D=\sequence{1,x,x^2,x^3} \)
      sebagai basis untuk kodomain.
      Dengan pilihan tersebut, kita memperoleh
      \begin{equation*}
        \rep{h}{B,D}
        =\begin{mat}[r]
           1   &1  &0  \\
           1   &2  &1  \\
           0   &0  &0  \\
           0   &0  &-1
         \end{mat}_{B,D}
      \end{equation*}
      dan karena
      \begin{equation*}
        \rep{1-3x+2x^2}{B}=\colvec[r]{1 \\ -3 \\ 2}_B
      \end{equation*}
      perhitungan perkalian matriks-vektor memberikan hasil berikut.
      \begin{equation*}
        \rep{h(1-3x+2x^2)}{D}=
         \begin{mat}[r]
           1   &1  &0  \\
           1   &2  &1  \\
           0   &0  &0  \\
           0   &0  &-1
         \end{mat}_{B,D}
         \colvec[r]{1 \\ -3 \\ 2}_B
         =\colvec[r]{-2 \\ -3 \\ 0 \\ -2}_D
      \end{equation*}
      Jadi, \( h(1-3x+2x^2)
              =-2\cdot 1-3\cdot x+0\cdot x^2-2\cdot x^3
              =-2-3x-2x^3 \),
      seperti di atas.  
    \end{answer}
  \item Misalkan $\map{h}{\Re^2}{\matspace_{\nbyn{2}}}$ adalah transformasi linear
   dengan aksi berikut.
   \begin{equation*}
     \colvec{1 \\ 0}\mapsto
     \begin{mat}
       1  &2  \\
       0  &1
     \end{mat}
     \qquad
     \colvec{0 \\ 1}\mapsto
     \begin{mat}
       0  &-1  \\
       1  &0
     \end{mat}
   \end{equation*}
   Apakah efeknya pada vektor umum dengan entri $x$ dan~$y$?
   \begin{answer}
     Tetapkan basis alami berikut untuk $\matspace_{\nbyn{2}}$.
     \begin{equation*}
       D=\sequence{
       \begin{mat}
          1 &0  \\
          0 &0
       \end{mat},
       \begin{mat}
          0 &1  \\
          0 &0
       \end{mat},
       \begin{mat}
          0 &0  \\
          1 &0
       \end{mat},
       \begin{mat}
          0 &0  \\
          0 &1
       \end{mat}}
     \end{equation*}
     Representasi peta~$h$ terhadap $\stdbasis_2,D$ adalah sebagai berikut.
     \begin{equation*}
       \rep{h}{\stdbasis_2,D}=
       \begin{mat}
         1  &0  \\
         2  &-1  \\
         0  &1  \\
         1  &0
       \end{mat}
     \end{equation*}
     Karena vektor umum direpresentasikan terhadap~$\stdbasis_2$ 
     oleh dirinya sendiri, 
     dan demikian pula setiap matriks direpresentasikan terhadap~$D$ oleh dirinya sendiri, 
     kita memperoleh efek peta berikut.
     \begin{equation*}
       \rep{h(\colvec{x \\ y})}{D}
       =
       \begin{mat}
         1  &0  \\
         2  &-1  \\
         0  &1  \\
         1  &0
       \end{mat}
       \colvec{x \\ y}
       =
       \colvec{x \\ 2x-y \\ y \\ x}
     \end{equation*}
   \end{answer}
  \recommended \item  
    Andaikan bahwa \( \map{h}{\Re^2}{\Re^3} \) ditentukan oleh 
    aksi berikut.
    \begin{equation*}
      \colvec[r]{1 \\ 0}\mapsto\colvec[r]{2 \\ 2 \\ 0}
      \qquad
      \colvec[r]{0 \\ 1}\mapsto\colvec[r]{0 \\ 1 \\ -1}
    \end{equation*}
    Dengan menggunakan basis standar, cari
    \begin{exparts}
      \partsitem matriks yang merepresentasikan peta ini;
      \partsitem rumus umum untuk \( h(\vec{v}) \).
    \end{exparts}
    \begin{answer}
      Sekali lagi, seperti diingatkan dalam subbagian ini, 
      terhadap $\stdbasis_i$, suatu vektor kolom merepresentasikan dirinya sendiri.
      \begin{exparts}
        \partsitem Untuk merepresentasikan \( h \) terhadap 
          \( \stdbasis_2,\stdbasis_3 \), ambil 
          bayangan vektor-vektor basis dari domain,
          lalu representasikan terhadap basis kodomain.
          Yang pertama adalah
          \begin{equation*}
            \rep{\,h(\vec{e}_1)\,}{\stdbasis_3}
            =\rep{\colvec[r]{2 \\ 2 \\ 0}}{\stdbasis_3}
            =\colvec[r]{2 \\ 2 \\ 0}
          \end{equation*}
          sedangkan yang kedua adalah berikut.
          \begin{equation*}
            \rep{\,h(\vec{e}_2)\,}{\stdbasis_3}
            =\rep{\colvec[r]{0 \\ 1 \\ -1}}{\stdbasis_3}
            =\colvec[r]{0 \\ 1 \\ -1}
          \end{equation*}
          Sejajarkan keduanya untuk membentuk matriks.
          \begin{equation*}
            \rep{h}{\stdbasis_2,\stdbasis_3}=
            \begin{mat}[r]
              2  &0  \\
              2  &1  \\
              0  &-1
            \end{mat}
          \end{equation*}
        \partsitem Untuk sebarang \( \vec{v} \) dalam domain \( \Re^2 \),
          \begin{equation*}
            \rep{\vec{v}}{\stdbasis_2}
            =\rep{\colvec{v_1 \\ v_2}}{\stdbasis_2}
            =\colvec{v_1 \\ v_2}
          \end{equation*}
          sehingga
          \begin{equation*}
            \rep{\,h(\vec{v})\,}{\stdbasis_3}
            =\begin{mat}[r]
              2  &0  \\
              2  &1  \\
              0  &-1
            \end{mat}
            \colvec{v_1 \\ v_2}
            =\colvec{2v_1 \\ 2v_1+v_2 \\ -v_2}
          \end{equation*}
          adalah representasi yang diinginkan.
      \end{exparts}  
    \end{answer}
  \item Representasikan homomorfisme $\map{h}{\Re^3}{\Re^2}$ yang diberikan oleh rumus ini
    terhadap basis-basis berikut.
    \begin{equation*}
      \colvec{x \\ y \\ z}\mapsto\colvec{x+y \\ x+z}
      \qquad
       B=\sequence{\colvec{1 \\ 1 \\ 1},
                   \colvec{1 \\ 1 \\ 0},
                   \colvec{1 \\ 0 \\ 0}}
       \quad
          D=\sequence{\colvec{1 \\ 0},
                  \colvec{0 \\ 2}}
    \end{equation*}
    \begin{answer}
      Aksi peta pada vektor-vektor basis domain adalah sebagai berikut.
      \begin{equation*}
          \colvec{1 \\ 1 \\ 1}\mapsto\colvec{2 \\ 2}
          \quad
          \colvec{1 \\ 1 \\ 0}\mapsto\colvec{2 \\ 1}
          \quad
          \colvec{1 \\ 0 \\ 0}\mapsto\colvec{1 \\ 1}
      \end{equation*}
      Representasikan semuanya terhadap basis kodomain.
      \begin{equation*}
        \rep{\colvec{2 \\ 2}}{D}=\colvec{2 \\ 1}_D
        \quad
        \rep{\colvec{2 \\ 1}}{D}=\colvec{2 \\ 1/2}_D
        \quad
        \rep{\colvec{1 \\ 1}}{D}=\colvec{1 \\ 1/2}_D
      \end{equation*}
      Konkatenasikan semuanya menjadi suatu matriks.
      \begin{equation*}
        \rep{h}{B,D}=
        \begin{mat}
          2 &2   &1 \\
          1 &1/2 &1/2
        \end{mat}
      \end{equation*}      
    \end{answer}
  \recommended \item  
    Misalkan \( \map{d/dx}{\polyspace_3}{\polyspace_3} \) adalah transformasi
    turunan.
    \begin{exparts}
      \partsitem Representasikan \( d/dx \) terhadap \( B,B \), dengan
        \( B=\sequence{1,x,x^2,x^3} \).
      \partsitem Representasikan \( d/dx \) terhadap \( B,D \), dengan
        \( D=\sequence{1,2x,3x^2,4x^3} \).
    \end{exparts}
    \begin{answer}
      \begin{exparts*}
        \partsitem 
        Mula-mula kita harus mencari bayangan setiap vektor dari basis domain,
        lalu merepresentasikan bayangan tersebut terhadap basis kodomain.
        \begin{multline*}
          \rep{\frac{d\,1}{dx}}{B}=\colvec[r]{0 \\ 0 \\ 0 \\ 0}
          \quad
          \rep{\frac{d\,x}{dx}}{B}=\colvec[r]{1 \\ 0 \\ 0 \\ 0}
          \quad
          \rep{\frac{d\,x^2}{dx}}{B}=\colvec[r]{0 \\ 2 \\ 0 \\ 0}
          \\
          \rep{\frac{d\,x^3}{dx}}{B}=\colvec[r]{0 \\ 0 \\ 3 \\ 0}
        \end{multline*}
        Representasi-representasi tersebut kemudian disejajarkan untuk membentuk matriks
        yang merepresentasikan peta.
        \begin{equation*}
          \rep{\frac{d}{dx}}{B,B}=
          \begin{mat}[r]
            0  &1  &0  &0  \\
            0  &0  &2  &0  \\
            0  &0  &0  &3  \\
            0  &0  &0  &0 
          \end{mat}  
        \end{equation*}
        \partsitem Dengan melanjutkan seperti pada butir sebelumnya, kita merepresentasikan bayangan
          vektor-vektor basis domain
        \begin{multline*}
          \rep{\frac{d\,1}{dx}}{D}=\colvec[r]{0 \\ 0 \\ 0 \\ 0}
          \quad
          \rep{\frac{d\,x}{dx}}{D}=\colvec[r]{1 \\ 0 \\ 0 \\ 0}
          \quad
          \rep{\frac{d\,x^2}{dx}}{D}=\colvec[r]{0 \\ 1 \\ 0 \\ 0}
          \\
          \rep{\frac{d\,x^3}{dx}}{D}=\colvec[r]{0 \\ 0 \\ 1 \\ 0}
        \end{multline*}
        lalu menyejajarkannya untuk membentuk matriks.
        \begin{equation*}
          \rep{\frac{d}{dx}}{B,D}=
          \begin{mat}[r]
            0  &1  &0  &0  \\
            0  &0  &1  &0  \\
            0  &0  &0  &1  \\
            0  &0  &0  &0
          \end{mat}  
        \end{equation*}
      \end{exparts*}  
    \end{answer}
  \recommended \item 
    Representasikan setiap peta linear terhadap setiap pasangan basis.
    \begin{exparts}
      \partsitem \( \map{d/dx}{\polyspace_n}{\polyspace_n} \) terhadap
        \( B,B \), dengan \( B=\sequence{1,x,\dots,x^n} \), yang diberikan oleh
        \begin{equation*}
            a_0+a_1x+a_2x^2+\dots+a_nx^n
            \mapsto
            a_1+2a_2x+\dots+na_nx^{n-1}
        \end{equation*}
      \partsitem \( \map{\int}{\polyspace_n}{\polyspace_{n+1}} \) 
        terhadap
        \( B_n,B_{n+1} \), dengan \( B_i=\sequence{1,x,\dots,x^i} \), yang diberikan oleh
        \begin{equation*}
            a_0+a_1x+a_2x^2+\dots+a_nx^n
            \mapsto
            a_0x+\frac{a_1}{2}x^2+\dots+\frac{a_n}{n+1}x^{n+1}
        \end{equation*}
      \partsitem \( \map{\int^1_0}{\polyspace_n}{\Re} \) terhadap
        \( B,\stdbasis_1 \), dengan \( B=\sequence{1,x,\dots,x^n} \)
        dan \( \stdbasis_1=\sequence{1} \), yang diberikan oleh
        \begin{equation*}
            a_0+a_1x+a_2x^2+\dots+a_nx^n
            \mapsto
            a_0+\frac{a_1}{2}+\dots+\frac{a_n}{n+1}
        \end{equation*}
      \partsitem \( \map{\text{eval}_3}{\polyspace_n}{\Re} \) terhadap
        \( B,\stdbasis_1 \), dengan \( B=\sequence{1,x,\dots,x^n} \)
        dan \( \stdbasis_1=\sequence{1} \), yang diberikan oleh
        \begin{equation*}
            a_0+a_1x+a_2x^2+\dots+a_nx^n
            \mapsto
            a_0+a_1\cdot 3+a_2\cdot 3^2+\dots+a_n\cdot 3^n
        \end{equation*}
      \partsitem \( \map{\text{slide}_{-1}}{\polyspace_n}{\polyspace_n} \) 
        terhadap
        \( B,B \), dengan \( B=\sequence{1,x,\ldots,x^n} \), yang diberikan oleh
        \begin{equation*}
            a_0+a_1x+a_2x^2+\dots+a_nx^n
            \mapsto
            a_0+a_1\cdot (x+1)+\dots+a_n\cdot (x+1)^n
        \end{equation*}
    \end{exparts}
    \begin{answer}
      Untuk masing-masing peta, kita harus mencari bayangan setiap vektor basis domain,
      merepresentasikan setiap bayangan terhadap basis kodomain,
      lalu menyejajarkan representasi-representasi tersebut untuk memperoleh matriks.
      \begin{exparts}
        \partsitem Vektor-vektor basis dari domain memiliki bayangan berikut
          \begin{equation*}
            1\mapsto 0  
            \quad x\mapsto 1  
            \quad x^2\mapsto 2x  
            \quad \ldots
          \end{equation*}
          dan bayangan-bayangan ini direpresentasikan terhadap basis
          kodomain dengan cara berikut.
          \begin{multline*}
            \rep{0}{B}=\colvec{0 \\ 0 \\ 0 \\ \vdots \\ \  \\  \ }
            \quad
            \rep{1}{B}=\colvec{1 \\ 0 \\ 0 \\ \vdots \\ \  \\ \ }
            \quad
            \rep{2x}{B}=\colvec{0 \\ 2 \\ 0 \\ \vdots \\ \  \\ \ }        \\
            \ldots\quad
            \rep{nx^{n-1}}{B}=\colvec{0 \\ 0 \\ 0 \\ \vdots \\ n \\ 0}
          \end{multline*}
          Matriks
          \begin{equation*}
            \rep{\frac{d}{dx}}{B,B}
            =\begin{mat}
              0  &1  &0  &\ldots  &0  \\
              0  &0  &2  &\ldots  &0  \\
                 &\vdots             \\
              0  &0  &0  &\ldots  &n  \\
              0  &0  &0  &\ldots  &0
            \end{mat}
          \end{equation*}
          memiliki $n+1$ baris dan kolom.
       \partsitem Setelah bayangan vektor-vektor basis domain di bawah peta ini
          ditentukan
          \begin{equation*}
            1\mapsto x 
            \quad x\mapsto x^2/2  
            \quad x^2\mapsto x^3/3 
            \quad \ldots
          \end{equation*}
          semuanya dapat direpresentasikan terhadap basis kodomain
          \begin{multline*}
            \rep{x}{B_{n+1}}=\colvec{0 \\ 1 \\ 0 \\ \vdots \\ \ }
            \quad
            \rep{x^2/2}{B_{n+1}}=\colvec{0 \\ 0 \\ 1/2 \\ \vdots \\ \ }   \\
            \ldots\quad
            \rep{x^{n+1}/(n+1)}{B_{n+1}}
                     =\colvec{0 \\ 0 \\ 0 \\ \vdots \\ 1/(n+1)}
          \end{multline*}
          lalu disatukan untuk membentuk matriks.
          \begin{equation*}
            \rep{\int}{B_{n},B_{n+1}}
            =\begin{mat}
              0  &0  &\ldots  &0      &0  \\
              1  &0  &\ldots  &0      &0  \\
              0  &1/2&\ldots  &0      &0  \\
                 &\vdots                  \\
              0  &0  &\ldots  &0      &1/(n+1)
            \end{mat}
          \end{equation*}
        \partsitem Bayangan vektor-vektor basis domain adalah
          \begin{equation*} 
            1\mapsto 1 
            \quad x\mapsto 1/2 
            \quad x^2\mapsto 1/3 
            \quad \ldots
          \end{equation*}
          dan semuanya direpresentasikan terhadap basis kodomain sebagai
          \begin{equation*}
            \rep{1}{\stdbasis_1}=1
            \quad \rep{1/2}{\stdbasis_1}=1/2
            \quad \ldots
          \end{equation*}
          sehingga matriksnya adalah 
          \begin{equation*}
            \rep{\int}{B,\stdbasis_1}
            =\begin{mat}
              1  &1/2 &\cdots  &1/n    &1/(n+1)
            \end{mat}
          \end{equation*}
          (ini adalah matriks $\nbym{1}{(n+1)}$).
        \partsitem Bayangan vektor-vektor basis domain adalah
          \begin{equation*}  
            1\mapsto 1 
            \quad x\mapsto 3 
            \quad x^2\mapsto 9
            \quad \ldots 
          \end{equation*}
          dan semuanya direpresentasikan dalam kodomain sebagai
          \begin{equation*}
            \rep{1}{\stdbasis_1}=1
            \quad\rep{3}{\stdbasis_1}=3
            \quad\rep{9}{\stdbasis_1}=9
            \quad \ldots 
          \end{equation*}
          sehingga matriksnya adalah sebagai berikut.
          \begin{equation*}
            \rep{\text{eval}_3}{B,\stdbasis_1}
            =\begin{mat}[r]
              1  &3   &9  &\cdots  &3^n
            \end{mat}
          \end{equation*}
        \partsitem Bayangan vektor-vektor basis dari domain adalah
          $1\mapsto 1$, 
          dan $x\mapsto x+1=1+x$,  
          dan $x^2\mapsto (x+1)^2=1+2x+x^2$,  
          dan $x^3\mapsto (x+1)^3=1+3x+3x^2+x^3$, dan seterusnya.  
          Representasi-representasinya adalah sebagai berikut.
          \begin{equation*}
            \rep{1}{B}=\colvec[r]{1 \\ 0 \\ 0 \\ 0 \\ \vdotswithin{0} \\ 0}
            \quad
            \rep{1+x}{B}=\colvec{1 \\ 1 \\ 0 \\ 0 \\ \vdotswithin{0} \\ 0}
            \quad
            \rep{1+2x+x^2}{B}=\colvec{1 \\ 2 \\ 1 \\ 0 \\ \vdotswithin{0} \\ 0}
            \quad\ldots
          \end{equation*}
          Matriks yang dihasilkan
          \begin{equation*}
            \renewcommand{\arraystretch}{1.2}
            \rep{\text{slide}_{-1}}{B,B}
            =\begin{mat}
              1  &1   &1  &1  &\ldots  &1           \\
              0  &1   &2  &3  &\ldots  &\binom{n}{1}  \\
              0  &0   &1  &3  &\ldots  &\binom{n}{2}  \\
                 &\vdots                        \\
              0  &0   &0  &   &\ldots  &1
            \end{mat}
          \end{equation*}
          adalah \definend{segitiga Pascal}\index{segitiga Pascal}
          (ingat bahwa $\binom{n}{r}$ adalah banyaknya cara memilih $r$
          objek, tanpa memperhatikan urutan dan tanpa pengulangan,
          dari suatu himpunan berukuran $n$). 
      \end{exparts}  
    \end{answer}
  \item 
    Representasikan peta identitas pada sebarang ruang
    taktrivial terhadap \( B,B \), dengan \( B \) sebarang basis.
    \begin{answer}
      Jika ruang tersebut berdimensi \( n \),
      \begin{equation*}
        \rep{\text{id}}{B,B}=
        \begin{mat}
          1  &0  \ldots  &0  \\
          0  &1  \ldots  &0  \\
             &\vdots         \\
          0  &0  \ldots  &1
        \end{mat}_{B,B}
      \end{equation*}
      adalah matriks identitas $\nbyn{n}$.  
    \end{answer}
  \item 
    Representasikan, terhadap basis alami, 
    transformasi transpos pada ruang 
    \( \matspace_{\nbyn{2}} \) matriks-matriks $\nbyn{2}$.
    \begin{answer}
      Ambil berikut ini sebagai basis alami
      \begin{equation*}
        B=\sequence{\vec{\beta}_1,\vec{\beta}_2,\vec{\beta}_3,\vec{\beta}_4}
         =\sequence{
            \begin{mat}[r]
              1  &0  \\
              0  &0
            \end{mat},
            \begin{mat}[r]
              0  &1  \\
              0  &0
            \end{mat},
            \begin{mat}[r]
              0  &0  \\
              1  &0
            \end{mat},
            \begin{mat}[r]
              0  &0  \\
              0  &1
            \end{mat}   }
      \end{equation*}
      maka peta transpos beraksi dengan cara berikut
      \begin{equation*}
        \vec{\beta}_1\mapsto\vec{\beta}_1 
        \quad \vec{\beta}_2\mapsto\vec{\beta}_3 
        \quad \vec{\beta}_3\mapsto\vec{\beta}_2 
        \quad \vec{\beta}_4\mapsto\vec{\beta}_4  
      \end{equation*}
      sehingga merepresentasikan bayangan-bayangan terhadap basis
      kodomain dan menyejajarkan vektor-vektor kolom tersebut memberikan hasil berikut.
      \begin{equation*}
        \rep{\text{trans}}{B,B}=
        \begin{mat}[r]
          1  &0  &0  &0  \\
          0  &0  &1  &0  \\
          0  &1  &0  &0  \\
          0  &0  &0  &1
        \end{mat}_{B,B}
      \end{equation*}   
     \end{answer}
  \item 
     Andaikan bahwa 
     \( B=\sequence{\vec{\beta}_1,\vec{\beta}_2,\vec{\beta}_3,\vec{\beta}_4} \)
     adalah basis suatu ruang vektor.
     Representasikan terhadap \( B,B \) transformasi yang ditentukan
     oleh masing-masing pengaitan berikut.
     \begin{exparts}
       \partsitem \( \vec{\beta}_1\mapsto\vec{\beta}_2 \),
         \( \vec{\beta}_2\mapsto\vec{\beta}_3 \),
         \( \vec{\beta}_3\mapsto\vec{\beta}_4 \),
         \( \vec{\beta}_4\mapsto\zero \)
       \partsitem \( \vec{\beta}_1\mapsto\vec{\beta}_2 \),
         \( \vec{\beta}_2\mapsto\zero \),
         \( \vec{\beta}_3\mapsto\vec{\beta}_4 \),
         \( \vec{\beta}_4\mapsto\zero \)
       \partsitem \( \vec{\beta}_1\mapsto\vec{\beta}_2 \),
         \( \vec{\beta}_2\mapsto\vec{\beta}_3 \),
         \( \vec{\beta}_3\mapsto\zero \),
         \( \vec{\beta}_4\mapsto\zero \)
     \end{exparts}
     \begin{answer}
       \begin{exparts}
         \partsitem Terhadap basis kodomain, bayangan anggota-anggota
           basis domain direpresentasikan sebagai
           \begin{equation*}
             \rep{\vec{\beta}_2}{B}=\colvec[r]{0 \\ 1 \\ 0 \\ 0}
             \quad
             \rep{\vec{\beta}_3}{B}=\colvec[r]{0 \\ 0 \\ 1 \\ 0}
             \quad
             \rep{\vec{\beta}_4}{B}=\colvec[r]{0 \\ 0 \\ 0 \\ 1}
             \quad
             \rep{\zero}{B}=\colvec[r]{0 \\ 0 \\ 0 \\ 0}
           \end{equation*}
           dan akibatnya, matriks yang merepresentasikan transformasi tersebut adalah 
           sebagai berikut.
           \begin{equation*}
             \begin{mat}[r]
                   0  &0  &0  &0  \\
                   1  &0  &0  &0  \\
                   0  &1  &0  &0  \\
                   0  &0  &1  &0
                 \end{mat}  
         \end{equation*}
         \partsitem 
              $\begin{mat}[r]
                   0  &0  &0  &0  \\
                   1  &0  &0  &0  \\
                   0  &0  &0  &0  \\
                   0  &0  &1  &0
                 \end{mat}$
         \partsitem 
               $\begin{mat}[r]
                   0  &0  &0  &0  \\
                   1  &0  &0  &0  \\
                   0  &1  &0  &0  \\
                   0  &0  &0  &0
                 \end{mat}$
       \end{exparts}  
     \end{answer}
  \item
    \nearbyexample{exam:RepsOfRigidPlaneMaps} menunjukkan cara merepresentasikan
    transformasi rotasi bidang terhadap 
    basis standar.
    Nyatakan pula transformasi-transformasi berikut terhadap basis
    standar.
    \begin{exparts}
      \partsitem peta \definend{dilatasi} $d_s$, yang mengalikan 
        semua vektor dengan skalar yang sama $s$\index{dilatasi!representasi}
      \partsitem peta \definend{refleksi} $f_\ell$, yang mencerminkan
        semua vektor terhadap garis $\ell$ yang melalui titik asal
    \end{exparts}
    \begin{answer}
      \begin{exparts}
        \partsitem Gambar $\map{d_s}{\Re^2}{\Re^2}$ adalah sebagai berikut.
          \begin{center}  \small
            \includegraphics{map/mp/ch3.75}
          \end{center}
         Efek peta ini pada vektor-vektor dalam basis standar domain 
         adalah 
         \begin{equation*} 
            \colvec[r]{1 \\ 0}\mapsunder{d_s}\colvec{s \\ 0}
            \qquad
            \colvec[r]{0 \\ 1}\mapsunder{d_s}\colvec{0 \\ s}
         \end{equation*}
         dan bayangan-bayangan tersebut direpresentasikan terhadap 
         basis kodomain (sekali lagi, basis standar) oleh dirinya sendiri.
         \begin{equation*} 
           \rep{\colvec{s \\ 0}}{\stdbasis_2}=\colvec{s \\ 0}
           \qquad
           \rep{\colvec{0 \\ s}}{\stdbasis_2}=\colvec{0 \\ s}
         \end{equation*}
         Jadi, representasi peta dilatasi adalah sebagai berikut.
         \begin{equation*}
           \rep{d_s}{\stdbasis_2,\stdbasis_2}
           =\begin{mat}
              s  &0  \\
              0  &s
           \end{mat}
        \end{equation*}
      \partsitem Gambar $\map{f_\ell}{\Re^2}{\Re^2}$ adalah sebagai berikut.
         \begin{center}  \small
           \includegraphics{map/mp/ch3.76}
        \end{center}
       Beberapa perhitungan (lihat Latihan~I.\ref{exer:RigidPlaneMapsAutos})
       menunjukkan bahwa ketika garis memiliki kemiringan $k$, 
       \begin{equation*}
         \colvec[r]{1 \\ 0}
             \mapsunder{f_\ell}\colvec{(1-k^2)/(1+k^2) \\ 2k/(1+k^2)}
         \qquad
         \colvec[r]{0 \\ 1}
             \mapsunder{f_\ell}\colvec{2k/(1+k^2) \\ -(1-k^2)/(1+k^2)}
       \end{equation*}
       (kasus garis dengan kemiringan takterdefinisi terpisah, tetapi mudah) 
       sehingga matriks yang merepresentasikan refleksi adalah sebagai berikut.
       \begin{equation*}
         \rep{f_\ell}{\stdbasis_2,\stdbasis_2}
         =\frac{1}{1+k^2}\cdot\begin{mat}
            1-k^2  &2k  \\
            2k     &-(1-k^2)
         \end{mat}
       \end{equation*}
      \end{exparts}  
    \end{answer}
  \recommended \item  
    Tinjau suatu transformasi linear pada \( \Re^2 \) yang ditentukan
    oleh kedua pengaitan berikut.
    \begin{equation*}
      \colvec[r]{1 \\ 1}\mapsto\colvec[r]{2 \\ 0}
      \qquad
      \colvec[r]{1 \\ 0}\mapsto\colvec[r]{-1 \\ 0}
    \end{equation*}
    \begin{exparts}
      \partsitem Representasikan transformasi ini terhadap basis-basis
        standar.
      \partsitem Ke manakah transformasi ini mengirim vektor berikut?
        \begin{equation*}
          \colvec[r]{0 \\ 5}
        \end{equation*}
      \partsitem Representasikan transformasi ini terhadap basis-basis berikut.
        \begin{equation*}
          B=\sequence{\colvec[r]{1 \\ -1},\colvec[r]{1 \\ 1}}
          \qquad
          D=\sequence{\colvec[r]{2 \\ 2},\colvec[r]{-1 \\ 1}}
        \end{equation*}
      \partsitem Dengan menggunakan \( B \) dari butir sebelumnya, 
        representasikan transformasi terhadap \( B,B \).
    \end{exparts}
    \begin{answer}
      Namai peta tersebut \( \map{t}{\Re^2}{\Re^2} \).
      \begin{exparts}
        \partsitem Untuk merepresentasikan peta ini terhadap basis-basis standar, kita
          harus mencari, lalu merepresentasikan, bayangan vektor $\vec{e}_1$
          dan $\vec{e}_2$ dari basis domain.
          Bayangan $\vec{e}_1$ telah diberikan.

          Salah satu cara mencari bayangan $\vec{e}_2$ adalah dengan 
          pengamatan langsung\Dash kita dapat melihat hal berikut.
          \begin{equation*}
            \colvec[r]{1 \\ 1}-\colvec[r]{1 \\ 0}=\colvec[r]{0 \\ 1}
            \;\mapsunder{t}\;
            \colvec[r]{2 \\ 0}-\colvec[r]{-1 \\ 0}=\colvec[r]{3 \\ 0}
          \end{equation*}

          Cara yang lebih sistematis untuk mencari bayangan $\vec{e}_2$ adalah
          menggunakan informasi yang diberikan untuk merepresentasikan transformasi, lalu 
          menggunakan representasi itu untuk menentukan bayangannya.  
          Dengan mengambil berikut ini sebagai basis,
          \begin{equation*}
            C=\sequence{\colvec[r]{1 \\ 1},\colvec[r]{1 \\ 0}}
          \end{equation*}
          informasi yang diberikan menyatakan hal berikut.
          \begin{equation*}
            \rep{t}{C,\stdbasis_2}=
            \begin{mat}[r]
              2  &-1  \\
              0  &0
            \end{mat}
          \end{equation*}
          Karena
          \begin{equation*}
            \rep{\vec{e}_2}{C}=\colvec[r]{1 \\ -1}_C
          \end{equation*}
kita memperoleh
          \begin{equation*}
            \rep{t(\vec{e}_2)}{\stdbasis_2}
            =\begin{mat}[r]
               2  &-1  \\
               0  &0
             \end{mat}_{C,\stdbasis_2}
            \colvec[r]{1 \\ -1}_C
            =\colvec[r]{3 \\ 0}_{\stdbasis_2}
          \end{equation*}
          dan akibatnya kita mengetahui bahwa $t(\vec{e}_2)=3\cdot\vec{e}_1$
          (karena terhadap basis standar, vektor ini 
          direpresentasikan oleh dirinya sendiri).
          Oleh karena itu, berikut adalah representasi $t$ terhadap
          $\stdbasis_2,\stdbasis_2$.
          \begin{equation*}
            \rep{t}{\stdbasis_2,\stdbasis_2}
            =\begin{mat}[r]
               -1 &3   \\
               0  &0
             \end{mat}_{\stdbasis_2,\stdbasis_2}
          \end{equation*}
        \partsitem Untuk menggunakan matriks yang diperoleh pada butir sebelumnya, perhatikan bahwa
          \begin{equation*}
            \rep{\colvec[r]{0 \\ 5}}{\stdbasis_2}=\colvec[r]{0 \\ 5}_{\stdbasis_2}
          \end{equation*}
          sehingga berikut adalah representasi bayangan vektor yang diberikan 
          terhadap basis kodomain.
          \begin{equation*}
            \rep{t(\colvec[r]{0 \\ 5})}{\stdbasis_2}
            =\begin{mat}[r]
              -1  &3   \\
               0  &0
             \end{mat}_{\stdbasis_2,\stdbasis_2}
            \colvec[r]{0 \\ 5}_{\stdbasis_2}
            =\colvec[r]{15 \\ 0}_{\stdbasis_2}
          \end{equation*}
          Karena basis kodomain adalah basis standar, sehingga vektor-vektor
          dalam kodomain direpresentasikan oleh dirinya sendiri, kita memperoleh hasil berikut.
          \begin{equation*}
            t(\colvec[r]{0 \\ 5})
            =\colvec[r]{15 \\ 0}
          \end{equation*}
        \partsitem Mula-mula kita mencari bayangan setiap anggota \( B \), lalu
          merepresentasikan bayangan-bayangan tersebut terhadap \( D \).
          Untuk langkah pertama, kita dapat menggunakan matriks yang diperoleh sebelumnya.
          \begin{equation*}
            \rep{\colvec[r]{1 \\-1}}{\stdbasis_2}
            =\begin{mat}[r]
              -1  &3   \\
               0  &0
             \end{mat}_{\stdbasis_2,\stdbasis_2}
            \colvec[r]{1 \\ -1}_{\stdbasis_2}
            =\colvec[r]{-4 \\ 0}_{\stdbasis_2}
            \quad\text{sehingga}\quad
            t(\colvec[r]{1 \\ -1})=\colvec[r]{-4 \\ 0}
          \end{equation*}
          Sebenarnya, untuk anggota kedua $B$ kita tidak perlu menerapkan
          matriks karena soal telah memberikan bayangannya.
          \begin{equation*}
            t(\colvec[r]{1 \\ 1})=\colvec[r]{2 \\ 0}
          \end{equation*}
          Kini merepresentasikan bayangan-bayangan tersebut terhadap $D$ merupakan pekerjaan rutin.
          \begin{equation*}
            \rep{\colvec[r]{-4 \\ 0}}{D}=\colvec[r]{-1 \\ 2}_D
            \quad\text{dan}\quad
            \rep{\colvec[r]{2 \\ 0}}{D}=\colvec[r]{1/2 \\ -1}_D
          \end{equation*}
          Jadi, matriksnya adalah sebagai berikut.
          \begin{equation*}
            \rep{t}{B,D}=
            \begin{mat}[r]
              -1  &1/2  \\
               2  &-1
            \end{mat}_{B,D}
          \end{equation*}
        \partsitem Kita mengetahui bayangan anggota-anggota basis domain
          dari butir sebelumnya.
          \begin{equation*}
             t(\colvec[r]{1 \\ -1})=\colvec[r]{-4 \\ 0}
             \qquad
             t(\colvec[r]{1 \\ 1})=\colvec[r]{2 \\ 0}
          \end{equation*}
          Kita dapat menghitung representasi bayangan-bayangan tersebut terhadap
          basis kodomain.
          \begin{equation*}
             \rep{\colvec[r]{-4 \\ 0}}{B}=\colvec[r]{-2 \\ -2}_B
             \quad\text{dan}\quad
             \rep{\colvec[r]{2 \\ 0}}{B}=\colvec[r]{1 \\ 1}_B
          \end{equation*}
          Jadi, berikut adalah matriksnya.
          \begin{equation*}
            \rep{t}{B,B}=
            \begin{mat}[r]
              -2  &1  \\
              -2  &1
            \end{mat}_{B,B}
          \end{equation*}
      \end{exparts}  
    \end{answer}
  \item 
    Andaikan bahwa \( \map{h}{V}{W} \) adalah satu-satu sehingga, 
    menurut \nearbytheorem{th:OOHomoEquivalence}, untuk sebarang basis
    \( B=\sequence{\vec{\beta}_1,\dots,\vec{\beta}_n}\subset V \), bayangan
    \( h(B)=\sequence{h(\vec{\beta}_1),\dots,h(\vec{\beta}_n)} \)
    merupakan basis untuk \( h(V) \).
    (Jika \( h \) adalah pada, maka \( h(V)=W \).)
    \begin{exparts}
      \partsitem Representasikan peta $h$ terhadap $\sequence{B,h(B)}$.
      \partsitem Untuk anggota $\vec{v}$ dari domain, dengan
        representasi $\vec{v}$ berkomponen $c_1$, \ldots, $c_n$,
        representasikan vektor bayangan \( h(\vec{v}) \) terhadap 
        basis bayangan $h(B)$.
    \end{exparts}
    \begin{answer}
      \begin{exparts}
        \partsitem Bayangan anggota-anggota basis domain adalah
          \begin{equation*}
            \vec{\beta}_1\mapsto h(\vec{\beta}_1)
            \quad 
            \vec{\beta}_2\mapsto h(\vec{\beta}_2)
            \quad\ldots\quad 
            \vec{\beta}_n\mapsto h(\vec{\beta}_n)
          \end{equation*}
          dan bayangan-bayangan tersebut direpresentasikan terhadap basis
          kodomain dengan cara berikut.
          \begin{multline*}
            \rep{\,h(\vec{\beta}_1)\,}{h(B)}=\colvec[r]{1 \\ 0 \\ \vdotswithin{0} \\ 0}
            \quad
            \rep{\,h(\vec{\beta}_2)\,}{h(B)}=\colvec[r]{0 \\ 1 \\ \vdotswithin{0} \\ 0}  \\
            \ldots\quad
            \rep{\,h(\vec{\beta}_n)\,}{h(B)}=\colvec[r]{0 \\ 0 \\ \vdotswithin{0} \\ 1}
          \end{multline*}
          Karena itu, matriksnya adalah matriks identitas.
          \begin{equation*}
            \rep{h}{B,h(B)}
            =\begin{mat}
               1  &0  &\ldots  &0  \\
               0  &1  &        &0  \\
                  &   &\ddots      \\
               0  &0  &        &1 
            \end{mat}
          \end{equation*}
        \partsitem Dengan menggunakan matriks pada butir sebelumnya, 
          representasinya adalah sebagai berikut.
          \begin{equation*}
            \rep{\,h(\vec{v})\,}{h(B)}
             =\colvec{c_1 \\ \vdots \\ c_n}_{h(B)}
          \end{equation*}
        \end{exparts}  
     \end{answer}
  \item 
    Berikan rumus untuk hasil kali suatu matriks dan \( \vec{e}_i \), yaitu
    vektor kolom yang semua entrinya nol kecuali satu entri bernilai satu pada
    posisi ke-\( i \).
    \begin{answer}
      Hasil kali
      \begin{equation*}
        \begin{mat}
          h_{1,1} &\ldots  &h_{1,i} &\ldots &h_{1,n}  \\
          h_{2,1} &\ldots  &h_{2,i} &\ldots &h_{2,n}  \\
                  &\vdots                             \\
          h_{m,1} &\ldots  &h_{m,i} &\ldots &h_{m,n}
        \end{mat}
        \colvec{0 \\ \vdots \\ 1 \\ \vdots \\ 0}
        =
        \colvec{h_{1,i} \\ h_{2,i} \\ \vdots \\ h_{m,i}}
      \end{equation*}
      memberikan kolom ke-\( i \) dari matriks.  
    \end{answer}
  \recommended \item 
    Untuk setiap ruang vektor fungsi satu variabel real,
    representasikan transformasi turunan terhadap \( B,B \).
    \begin{exparts}
      \partsitem \( \set{a\cos x+b\sin x \suchthat a,b\in\Re} \),
         \( B=\sequence{\cos x,\sin x} \)
      \partsitem \( \set{ae^x+be^{2x} \suchthat a,b\in\Re} \),
         \( B=\sequence{e^x,e^{2x}} \)
      \partsitem \( \set{a+bx+ce^x+dxe^{x} \suchthat a,b,c,d\in\Re} \),
         \( B=\sequence{1,x,e^x,xe^{x}} \)
    \end{exparts}
    \begin{answer}
      \begin{exparts}
        \partsitem Bayangan vektor-vektor basis domain adalah
         \( \cos x\mapsunder{d/dx}-\sin x \) dan
          \( \sin x\mapsunder{d/dx}\cos x \).
          Merepresentasikannya terhadap basis kodomain (sekali lagi, $B$)
          dan menyejajarkan representasi-representasi tersebut memberikan matriks berikut.
          \begin{equation*}
            \rep{\frac{d}{dx}}{B,B}=
            \begin{mat}[r]
                0  &1  \\
               -1  &0
            \end{mat}_{B,B}
          \end{equation*}
        \partsitem Bayangan vektor-vektor dalam basis domain adalah 
          \( e^x\mapsunder{d/dx}e^x \) dan
          \( e^{2x}\mapsunder{d/dx}2e^{2x} \).
          Merepresentasikannya terhadap basis kodomain dan menyejajarkannya
          memberikan matriks berikut.
          \begin{equation*}
            \rep{\frac{d}{dx}}{B,B}=
            \begin{mat}[r]
                1  &0  \\
                0  &2
            \end{mat}_{B,B}
          \end{equation*}
        \partsitem Bayangan anggota-anggota basis domain adalah 
          \( 1\mapsunder{d/dx}0 \), 
          \( x\mapsunder{d/dx}1 \),
          \( e^{x}\mapsunder{d/dx}e^{x} \), dan
          \( xe^{x}\mapsunder{d/dx}e^x+xe^x \).
          Merepresentasikan bayangan-bayangan ini terhadap $B$ dan menyejajarkannya
          memberikan matriks berikut.
          \begin{equation*}
            \rep{\frac{d}{dx}}{B,B}=
            \begin{mat}[r]
                0  &1  &0  &0 \\
                0  &0  &0  &0 \\
                0  &0  &1  &1 \\
                0  &0  &0  &1
            \end{mat}_{B,B}
          \end{equation*}
      \end{exparts}  
    \end{answer}
  \item 
    Cari jangkauan transformasi linear pada \( \Re^2 \) yang direpresentasikan
    terhadap basis-basis standar oleh setiap matriks berikut.
    \begin{exparts*}
      \partsitem $\begin{mat}[r]
            1  &0 \\
            0  &0
          \end{mat}$
      \partsitem $\begin{mat}[r]
          0  &0  \\
          3  &2  
        \end{mat}$
      \partsitem matriks berbentuk 
        $\begin{mat}
            a   &b  \\
            2a  &2b
          \end{mat}$
    \end{exparts*}
    \begin{answer}
      \begin{exparts}
        \partsitem Jangkauannya adalah himpunan vektor kodomain yang direpresentasikan
          terhadap basis kodomain dengan cara berikut.
          \begin{equation*}
            \set{
              \begin{mat}[r]
                1  &0  \\
                0  &0
              \end{mat}
              \colvec{x  \\ y}
              \suchthat x,y\in\Re}
            =\set{\colvec{x  \\ 0}
                  \suchthat x,y\in\Re}
          \end{equation*}
          Karena basis kodomain adalah $\stdbasis_2$, 
          sehingga setiap vektor direpresentasikan
          oleh dirinya sendiri, jangkauan transformasi ini adalah sumbu-$x$.
        \partsitem Jangkauannya adalah himpunan vektor kodomain yang direpresentasikan
          dengan cara berikut.
          \begin{equation*}
            \set{
              \begin{mat}[r]
                0  &0  \\
                3  &2
              \end{mat}
              \colvec{x  \\ y}
              \suchthat x,y\in\Re}
            =\set{\colvec{0  \\ 3x+2y}
                  \suchthat x,y\in\Re}
          \end{equation*}
          Terhadap $\stdbasis_2$, vektor-vektor merepresentasikan 
          dirinya sendiri, sehingga jangkauan ini
          adalah sumbu-$y$.
        \partsitem Himpunan vektor yang direpresentasikan 
          terhadap $\stdbasis_2$ sebagai
          \begin{align*}
            \set{
              \begin{mat}
                a   &b  \\
                2a  &2b
              \end{mat}
              \colvec{x  \\ y}
              \suchthat x,y\in\Re}
            &=\set{\colvec{ax+by  \\ 2ax+2by}
                  \suchthat x,y\in\Re}                    \\
            &=\set{(ax+by)\cdot\colvec[r]{1  \\ 2}
                  \suchthat x,y\in\Re}
          \end{align*}
          adalah garis $y=2x$, asalkan $a$ atau $b$ tidak nol, dan
          merupakan himpunan yang hanya terdiri atas titik asal jika keduanya nol.
      \end{exparts}  
    \end{answer}
  \recommended \item  
    Dapatkah satu matriks merepresentasikan dua peta linear yang berbeda?
    Dengan kata lain, dapatkah \( \rep{h}{B,D}=\rep{\hat{h}}{\hat{B},\hat{D}} \)?
    \begin{answer}
      Ya, karena dua alasan.

      Pertama, kedua peta $h$ dan $\hat{h}$ tidak harus memiliki domain
      dan kodomain yang sama.
      Sebagai contoh,
      \begin{equation*}
        \begin{mat}[r]
          1  &2  \\
          3  &4
        \end{mat}
      \end{equation*}
      merepresentasikan peta \( \map{h}{\Re^2}{\Re^2} \) terhadap basis-basis
      standar yang memetakan
      \begin{equation*}
        \colvec[r]{1 \\ 0}\mapsto\colvec[r]{1 \\ 3}
        \quad\text{dan}\quad
        \colvec[r]{0 \\ 1}\mapsto\colvec[r]{2 \\ 4}
      \end{equation*}
      serta juga merepresentasikan suatu
      \( \map{\hat{h}}{\polyspace_1}{\Re^2} \) terhadap
      \( \sequence{1,x} \) dan \( \stdbasis_2 \) yang beraksi dengan cara berikut.
      \begin{equation*}
        1\mapsto\colvec[r]{1 \\ 3}
        \quad\text{dan}\quad
        x\mapsto\colvec[r]{2 \\ 4}
      \end{equation*}

      Alasan kedua adalah bahwa, sekalipun domain dan
      kodomain \( h \) dan \( \hat{h} \) sama, basis yang berbeda menghasilkan
      peta yang berbeda.
      Contohnya adalah matriks identitas $\nbyn{2}$
      \begin{equation*}
        I=\begin{mat}[r]
          1  &0  \\
          0  &1
        \end{mat}
      \end{equation*}
      yang merepresentasikan peta identitas pada $\Re^2$ terhadap
      $\stdbasis_2,\stdbasis_2$.
      Namun, terhadap $\stdbasis_2$ untuk domain tetapi basis 
      $D=\sequence{\vec{e}_2,\vec{e}_1}$ untuk kodomain,
      matriks $I$ yang sama merepresentasikan peta yang menukar komponen pertama dan
      kedua 
      \begin{equation*}
        \colvec{x \\ y}\mapsto\colvec{y \\ x}
      \end{equation*}
      (yakni refleksi terhadap garis $y=x$).
    \end{answer}
  \item \label{exer:MatVecMultRepLinMap} 
    Buktikan \nearbytheorem{th:MatMultRepsFuncAppl}.
    \begin{answer}
      Kita meniru \nearbyexample{ex:TypLinMapRepByMat}, hanya dengan mengganti 
      bilangan dengan huruf.

      Tuliskan \( B \) sebagai \( \sequence{\vec{\beta}_1,\ldots,\vec{\beta}_n} \)
      dan \( D \) sebagai \( \sequence{\vec{\delta}_1,\dots,\vec{\delta}_m} \).
      Menurut definisi representasi suatu peta terhadap basis-basis,
      asumsi bahwa
      \begin{equation*}
        \rep{h}{B,D}
        =\begin{mat}
           h_{1,1} &\ldots  &h_{1,n}  \\
           \vdots  &        &\vdots   \\
           h_{m,1} &\ldots  &h_{m,n}
         \end{mat}
      \end{equation*}
      berarti bahwa
      $h(\vec{\beta}_i)=h_{1,i}\vec{\delta}_1+\dots+h_{m,i}\vec{\delta}_m$.
      Dan, menurut definisi representasi suatu vektor terhadap
      suatu basis, asumsi bahwa
      \begin{equation*}
        \rep{\vec{v}}{B}=\colvec{c_1 \\ \vdots \\ c_n}
      \end{equation*}
      berarti bahwa \( \vec{v}=c_1\vec{\beta}_1+\cdots+c_n\vec{\beta}_n \).
      Substitusi memberikan
      \begin{align*}
        h(\vec{v})
        &=h(c_1\cdot\vec{\beta}_1+\dots+c_n\cdot\vec{\beta}_n)      \\
        &=c_1\cdot h(\vec{\beta}_1)+\dots+c_n\cdot h(\vec{\beta}_n)    \\
        &=c_1\cdot (h_{1,1}\vec{\delta}_1+\dots+h_{m,1}\vec{\delta}_m) 
        +\dots                                         
        +c_n\cdot (h_{1,n}\vec{\delta}_1+\dots+h_{m,n}\vec{\delta}_m) \\
        &=(h_{1,1}c_1+\dots+h_{1,n}c_n)\cdot\vec{\delta}_1    
        +\cdots                                
        +(h_{m,1}c_1+\dots+h_{m,n}c_n)\cdot\vec{\delta}_m
      \end{align*}
      sehingga $h(\vec{v})$ direpresentasikan seperti yang diperlukan.   
    \end{answer}
  \recommended \item  
    \nearbyexample{exam:RepsOfRigidPlaneMaps} menunjukkan cara merepresentasikan 
    rotasi semua vektor di bidang sebesar sudut
    \( \theta \) terhadap titik asal,
    terhadap basis-basis standar.
    \begin{exparts}
      \partsitem Rotasi semua vektor dalam ruang-dimensi-tiga sebesar sudut
        \( \theta \) terhadap sumbu-\( x \) adalah suatu transformasi pada $\Re^3$.
        Representasikan terhadap basis-basis standar.
        Atur rotasi sedemikian rupa sehingga 
        bagi seseorang yang kakinya berada di titik asal dan
        kepalanya berada di \( (1,0,0) \), gerakannya tampak searah jarum jam.
      \partsitem Ulangi butir sebelumnya, tetapi kali ini rotasikan terhadap sumbu-\( y \).
        (Letakkan kepala orang tersebut di $\vec{e}_2$.)
      \partsitem Ulangi terhadap sumbu-\( z \).
      \partsitem Perluas butir sebelumnya ke \( \Re^4 \).
        (\textit{Petunjuk:} kita dapat menyatakan kembali 
        `rotasi terhadap sumbu-\( z \)' sebagai `rotasi sejajar
        dengan bidang \( xy \)'.)
    \end{exparts}
    \begin{answer}
      \begin{exparts}
        \partsitem Gambarnya adalah sebagai berikut.
          \begin{center}  \small
            \includegraphics{map/mp/ch3.77}
         \end{center}
         Bayangan vektor-vektor dari basis domain 
         \begin{equation*}
           \colvec[r]{1 \\ 0 \\ 0}\mapsto\colvec[r]{1 \\ 0 \\ 0}
           \qquad
           \colvec[r]{0 \\ 1 \\ 0}\mapsto\colvec{0 \\ \cos\theta \\ -\sin\theta}
           \qquad
           \colvec[r]{0 \\ 0 \\ 1}\mapsto\colvec{0 \\ \sin\theta \\ \cos\theta}
         \end{equation*}
         direpresentasikan terhadap basis kodomain
         (sekali lagi, $\stdbasis_3$) oleh dirinya sendiri, 
         sehingga menyejajarkan representasi-representasi tersebut untuk
         membentuk matriks memberikan hasil berikut.
         \begin{equation*}
            \rep{r_\theta}{\stdbasis_3,\stdbasis_3}=
            \begin{mat}
              1  &0       &0                \\
              0  &\cos\theta   &\sin\theta   \\
              0  &-\sin\theta  &\cos\theta
            \end{mat}                  
         \end{equation*}
       \partsitem Gambarnya serupa dengan gambar pada jawaban sebelumnya. 
         Bayangan vektor-vektor dari basis domain 
         \begin{equation*}
           \colvec[r]{1 \\ 0 \\ 0}\mapsto\colvec{\cos\theta \\ 0 \\ \sin\theta}
           \qquad
           \colvec[r]{0 \\ 1 \\ 0}\mapsto\colvec[r]{0 \\ 1 \\ 0}
           \qquad
           \colvec[r]{0 \\ 0 \\ 1}\mapsto\colvec{-\sin\theta \\ 0 \\ \cos\theta}
         \end{equation*}
         direpresentasikan terhadap basis kodomain $\stdbasis_3$
         oleh dirinya sendiri, sehingga berikut adalah matriksnya.
         \begin{equation*}
            \begin{mat}
              \cos\theta  &0       &-\sin\theta   \\
              0           &1       &0             \\
              \sin\theta  &0       &\cos\theta
            \end{mat}                  
         \end{equation*}
       \partsitem Bagi seseorang yang berdiri tegak, dengan sumbu-$z$ vertikal,
         rotasi bidang $xy$ yang searah jarum jam bergerak dari
         sumbu-$y$ positif menuju sumbu-$x$ positif.
         Artinya, rotasi tersebut berlawanan dengan arah dalam  
         \nearbyexample{exam:RepsOfRigidPlaneMaps}.
         Bayangan vektor-vektor dari basis domain 
         \begin{equation*}
           \colvec[r]{1 \\ 0 \\ 0}\mapsto\colvec{\cos\theta \\ -\sin\theta \\ 0}
           \qquad
           \colvec[r]{0 \\ 1 \\ 0}\mapsto\colvec{\sin\theta \\ \cos\theta \\ 0}
           \qquad
           \colvec[r]{0 \\ 0 \\ 1}\mapsto\colvec[r]{0 \\ 0 \\ 1}
         \end{equation*}
         direpresentasikan terhadap $\stdbasis_3$
         oleh dirinya sendiri, sehingga matriksnya adalah sebagai berikut.
         \begin{equation*}
            \begin{mat}
              \cos\theta    &\sin\theta  &0   \\
              -\sin\theta   &\cos\theta  &0    \\
              0             &0           &1
            \end{mat}                  
         \end{equation*}
        \partsitem 
            $\begin{mat}
              \cos\theta  &\sin\theta &0 &0 \\
              -\sin\theta  &\cos\theta  &0 &0 \\
              0           &0           &1 &0 \\
              0           &0           &0 &1
            \end{mat}$
      \end{exparts}   
    \end{answer}
  \item (Lema Triangularisasi Schur)\index{triangularisasi}%
    \index{matriks!segitiga}
    \begin{exparts}
      \partsitem Misalkan \( U \) adalah subruang dari \( V \) dan tetapkan basis-basis
        \( B_U\subseteq B_V \).
        Apakah hubungan antara representasi suatu vektor
        dari \( U \) terhadap
        \( B_U \) dan representasi vektor tersebut
        (dipandang sebagai anggota \( V \)) terhadap
        \( B_V \)?
      \partsitem Bagaimana dengan peta?
      \partsitem Tetapkan suatu basis 
        \( B=\sequence{\vec{\beta}_1,\dots,\vec{\beta}_n} \)
        untuk \( V \) dan amati bahwa rentang-rentang
        \begin{equation*}
          \spanof{\emptyset}=\set{\zero}\subset\spanof{\set{\vec{\beta}_1}}
                        \subset\spanof{\set{\vec{\beta}_1,\vec{\beta}_2}}
               \subset \quad\cdots\quad 
               \subset\spanof{B}=V
        \end{equation*}
        membentuk suatu rantai subruang yang meningkat secara ketat.
        Tunjukkan bahwa untuk sebarang peta linear \( \map{h}{V}{W} \), terdapat rantai
        \( W_0=\set{\zero}\subseteq W_1\subseteq \dots \subseteq W_m =W \)
        subruang dari \( W \) sedemikian sehingga 
        \begin{equation*}
          h(\spanof{\set{\vec{\beta}_1,\dots,\vec{\beta}_i}})\subseteq W_i
        \end{equation*}
        untuk setiap \( i \).
      \partsitem Simpulkan bahwa untuk setiap peta linear 
        \( \map{h}{V}{W} \), terdapat
        basis \( B,D \) sehingga matriks yang merepresentasikan \( h \) terhadap
        \( B,D \) berbentuk segitiga atas
        (artinya, setiap entri \( h_{i,j} \) dengan \( i>j \) adalah nol).
      \item Apakah representasi segitiga atas itu unik?
    \end{exparts}
    \begin{answer}
      \begin{exparts}
        \partsitem 
          Tuliskan basis \( B_U \) sebagai 
          \( \sequence{\vec{\beta}_1,\dots,\vec{\beta}_k} \) dan
          lalu tuliskan $B_V$ sebagai perluasan 
          \( \sequence{\vec{\beta}_1,\dots,\vec{\beta}_k,
                            \vec{\beta}_{k+1},\dots,\vec{\beta}_n} \).
          Jika 
          \begin{equation*}
            \rep{\vec{v}}{B_U}=\colvec{c_1 \\ \vdots \\ c_k}
          \end{equation*}
          sehingga $\vec{v}=c_1\cdot\vec{\beta}_1+\cdots+c_k\cdot\vec{\beta}_k$,
          maka
          \begin{equation*}
            \rep{\vec{v}}{B_V}=\colvec{c_1 \\ \vdots\\ c_k \\ 0 \\ \vdots \\ 0}
          \end{equation*}
          karena $\vec{v}=c_1\cdot\vec{\beta}_1+\dots+c_k\cdot\vec{\beta}_k
                    +0\cdot\vec{\beta}_{k+1}+\dots+0\cdot\vec{\beta}_n$.
        \partsitem
          Mula-mula kita harus menentukan arti pertanyaan tersebut.
          Bandingkan \( \map{h}{V}{W} \) dengan pembatasannya pada subruang
          \( \map{\restrictionmap{h}{U}}{U}{W} \).
          Ruang jangkauan pembatasan tersebut merupakan subruang dari \( W \), jadi tetapkan
          basis \( D_{h(U)} \) untuk ruang jangkauan ini dan perluas menjadi basis 
          \( D_V \) untuk \( W \).
          Kita mencari hubungan antara kedua hal berikut.
          \begin{equation*}
            \rep{h}{B_V,D_V}
            \quad\text{dan}\quad
            \rep{\restrictionmap{h}{U}}{B_U,D_{h(U)}}
          \end{equation*}
          Jawabannya langsung mengikuti butir sebelumnya:~jika
          \begin{equation*}
            \rep{\restrictionmap{h}{U}}{B_U,D_{h(U)}}
             =\begin{mat}
                h_{1,1}  &\ldots  &h_{1,k}  \\
                \vdots   &        &\vdots   \\
                h_{p,1}  &\ldots  &h_{p,k}
              \end{mat}
          \end{equation*}
          maka perluasannya direpresentasikan dengan cara berikut.
          \begin{equation*}
            \rep{h}{B_V,D_V}
             =\begin{mat}
                h_{1,1}  &\ldots  &h_{1,k}  &h_{1,k+1}  &\ldots  &h_{1,n}  \\
                \vdots   &        &         &           &        &\vdots   \\
                h_{p,1}  &\ldots  &h_{p,k}  &h_{p,k+1}  &\ldots  &h_{p,n}  \\
                0        &\ldots  &0        &h_{p+1,k+1}&\ldots  &h_{p+1,n}  \\
                \vdots   &        &         &           &        &\vdots   \\
                0        &\ldots  &0        &h_{m,k+1}  &\ldots  &h_{m,n}
              \end{mat}
          \end{equation*}
        \partsitem Ambil \( W_i \) sebagai rentang dari
          \( \set{h(\vec{\beta}_1),\dots,h(\vec{\beta}_i)} \).
        \partsitem Terapkan jawaban butir kedua pada butir ketiga.
        \partsitem Tidak.
          Sebagai contoh, \( \map{\pi_x}{\Re^2}{\Re^2} \), proyeksi pada
          sumbu-\( x \), direpresentasikan oleh kedua matriks segitiga atas
          berikut 
          \begin{equation*}
             \rep{\pi_x}{\stdbasis_2,\stdbasis_2}=
             \begin{mat}[r]
               1  &0  \\
               0  &0
             \end{mat}
             \quad\text{dan}\quad
             \rep{\pi_x}{C,\stdbasis_2}=
             \begin{mat}[r]
               0  &1  \\
               0  &0
             \end{mat}
          \end{equation*}
          dengan \( C=\sequence{\vec{e}_2,\vec{e}_1} \).
      \end{exparts}  
    \end{answer}
\index{homomorfisme!matriks yang merepresentasikan|)}
\end{exercises}















\subsection{Setiap Matriks Merepresentasikan Peta Linear}

%<*MapDefinedByMat0>
Subbagian sebelumnya menunjukkan bahwa
aksi suatu peta linear $h$ dijelaskan oleh matriks $H$,
terhadap basis-basis yang sesuai, dengan cara berikut.
\begin{equation*}
 \vec{v}=\colvec{v_1 \\ \vdots \\ v_n}_B
  \quad\overset{h}{\underset{H}{\longmapsto}}\quad
  h(\vec{v})=
  \colvec{h_{1,1}v_1+\dots+h_{1,n}v_n \\ 
                     \vdots                      \\
                     h_{m,1}v_1+\dots+h_{m,n}v_n}_D
  \tag{$*$}
\end{equation*}
Di sini kita akan menunjukkan kebalikannya, yakni bahwa
setiap matriks merepresentasikan suatu peta linear.
%</MapDefinedByMat0>

%<*MapDefinedByMat1>
Jadi, kita mulai dengan suatu matriks
\begin{equation*}
  H=\generalmatrix{h}{n}{m}
\end{equation*}
dan akan menjelaskan cara matriks tersebut mendefinisikan suatu peta~$h$.
Kita mengharuskan peta itu direpresentasikan oleh matriks, jadi 
mula-mula perhatikan bahwa dalam ($*$), dimensi domain peta adalah
banyaknya kolom~$n$ dari matriks 
dan dimensi kodomain adalah banyaknya baris~$m$.
Karena itu, sebagai domain $h$, 
tetapkan suatu ruang vektor $V$ berdimensi \( n \), dan sebagai
kodomain tetapkan suatu ruang $W$ berdimensi \( m \).
Tetapkan pula basis
\( B=\sequence{\vec{\beta}_1,\dots,\vec{\beta}_n} \) dan
\( D=\sequence{\vec{\delta}_1,\dots,\vec{\delta}_m} \) untuk ruang-ruang tersebut.
%</MapDefinedByMat1>

%<*MapDefinedByMat2>
Sekarang definisikan \( \map{h}{V}{W} \) sebagai berikut:~jika $\vec{v}$ dalam domain
memiliki representasi
\begin{equation*}
  \rep{\vec{v}}{B}
    =\colvec{v_1 \\ \vdots \\ v_n}_B
\end{equation*}
maka bayangannya \( h(\vec{v}) \) adalah anggota kodomain 
dengan representasi berikut.
\begin{equation*}
  \rep{\,h(\vec{v})\,}{D}
    =\colvec{h_{1,1}v_1+\dots+h_{1,n}v_n \\ \vdots \\ 
                   h_{m,1}v_1+\dots+h_{m,n}v_n}_D
\end{equation*}
Dengan kata lain, untuk menghitung aksi $h$ pada 
sebarang $\vec{v}\in V$, mula-mula nyatakan $\vec{v}$ terhadap
basis sebagai
$\vec{v}=v_1\vec{\beta}_1+\dots+v_n\vec{\beta}_n$, lalu 
$h(\vec{v})=
 (h_{1,1}v_1+\dots+h_{1,n}v_n)\cdot\vec{\delta}_1
  +\dots+
  (h_{m,1}v_1+\dots+h_{m,n}v_n)\cdot\vec{\delta}_m$.
%</MapDefinedByMat2>

%<*MapDefinedByMat3>
Di atas kita telah membuat beberapa pilihan; misalnya, $V$ dapat berupa sebarang
ruang berdimensi-$n$ dan $B$ dapat berupa sebarang basis untuk $V$,
jadi $H$ tidak mendefinisikan fungsi yang unik.
Namun, perhatikan bahwa   
setelah kita menetapkan $V$, $B$, $W$, dan $D$, 
$h$ terdefinisi dengan baik
karena $\vec{v}$ memiliki representasi unik terhadap basis $B$
% by Theorem~II.\ref{th:BasisIffUniqueRepWRT} 
dan penentuan 
$\vec{w}$ dari representasinya juga bersifat unik.
%</MapDefinedByMat3>

\begin{example}
Tinjau matriks berikut.
\begin{equation*}
  H=
  \begin{mat}
    1 &2 \\
    3 &4 \\
    5 &6
  \end{mat}
\end{equation*}
Matriks ini berukuran $\nbym{3}{2}$, sehingga setiap peta yang didefinisikannya harus membawa domain
berdimensi~$2$ ke kodomain berdimensi~$3$.
Kita dapat memilih domain dan kodomain sebagai $\Re^2$ dan $\polyspace_2$, 
dengan basis-basis berikut.
\begin{equation*}
  B=\sequence{\colvec{1 \\ 1}, \colvec[r]{1 \\ -1}}
  \qquad
  D=\sequence{x^2, x^2+x, x^2+x+1}
\end{equation*}
Kemudian, misalkan $\map{h}{\Re^2}{\polyspace_2}$ adalah fungsi yang didefinisikan oleh~$H$.
Kita akan menghitung bayangan anggota domain berikut di bawah $h$.
\begin{equation*}
  \vec{v}=\colvec[r]{-3 \\ 2}
\end{equation*}
Perhitungannya langsung.
\begin{equation*}
  \rep{h(\vec{v})}{D}
  =H\cdot\rep{\vec{v}}{B}
  =
  \begin{mat}
    1  &2  \\
    3  &4  \\
    5 &6
  \end{mat}
  \colvec{-1/2 \\ -5/2}
  =\colvec{-11/2 \\ -23/2 \\ -35/2}
\end{equation*}
Dari representasinya, perhitungan $h(\vec{v})$ merupakan pekerjaan rutin:
$(-11/2)(x^2)-(23/2)(x^2+x)-(35/2)(x^2+x+1)
=(-69/2)x^2-(58/2)x-(35/2)$.  
\end{example}


\begin{theorem} \label{th:MatIsLinMap}\index{homomorfisme!matriks yang merepresentasikan}
%<*th:MatIsLinMap>
Setiap matriks merepresentasikan suatu homomorfisme antara ruang-ruang vektor
berdimensi sesuai, terhadap sebarang pasangan basis.
%</th:MatIsLinMap>
\end{theorem}

\begin{proof}
%<*pf:MatIsLinMap>
Kita harus memeriksa bahwa untuk sebarang matriks $H$ dan sebarang basis domain serta kodomain $B,D$,
peta \( h \) yang didefinisikan bersifat linear.
Jika $\vec{v}, \vec{u}\in V$ sedemikian sehingga
\begin{equation*}
  \rep{\vec{v}}{B}=\colvec{v_1 \\ \vdots \\ v_n}
    \qquad
  \rep{\vec{u}}{B}=\colvec{u_1 \\ \vdots \\ u_n}
\end{equation*}
dan \( c,d\in\Re \), maka perhitungan
\begin{align*}
  h(c\vec{v}+d\vec{u})
  &=\bigl(h_{1,1}(cv_1+du_1)+\dots+
          h_{1,n}(cv_n+du_n)\bigr)\cdot\vec{\delta}_1+  \\
  & \hbox{}\quad\cdots+\bigl(h_{m,1}(cv_1+du_1)+\dots
         +h_{m,n}(cv_n+du_n)\bigr)\cdot\vec{\delta}_m  \\
  &=c\cdot h(\vec{v})+d\cdot h(\vec{u})
\end{align*}
memberikan pemeriksaan tersebut.
%</pf:MatIsLinMap>
\end{proof}

\begin{example} \label{ex:CngBasesChgMap}
Sekalipun domain dan kodomainnya sama, peta yang 
direpresentasikan matriks bergantung pada 
basis yang kita pilih.
Jika
\begin{equation*}
  H=
  \begin{mat}[r]
    1  &0  \\
    0  &0
  \end{mat},
  \quad
  B_1=D_1=\sequence{\colvec[r]{1 \\ 0},\colvec[r]{0 \\ 1} },
  \quad\text{dan}\quad
  B_2=D_2=\sequence{\colvec[r]{0 \\ 1},\colvec[r]{1 \\ 0} },
\end{equation*}
maka \( \map{h_1}{\Re^2}{\Re^2} \) yang direpresentasikan oleh \( H \)
terhadap \( B_1,D_1 \) memetakan
\begin{equation*}
  \colvec{c_1 \\ c_2}
  =\colvec{c_1 \\ c_2}_{B_1}
  \quad
  \mapsto
  \quad
  \colvec{c_1 \\ 0}_{D_1}
  =
  \colvec{c_1 \\ 0}
\end{equation*}
sedangkan \( \map{h_2}{\Re^2}{\Re^2} \) yang direpresentasikan oleh \( H \)
terhadap \( B_2,D_2 \) adalah peta berikut.
\begin{equation*}
  \colvec{c_1 \\ c_2}
  =\colvec{c_2 \\ c_1}_{B_2}
  \quad
  \mapsto
  \quad
  \colvec{c_2 \\ 0}_{D_2}
  =
  \colvec{0 \\ c_2}
\end{equation*}
Keduanya adalah fungsi yang berbeda.
Yang pertama merupakan proyeksi pada sumbu-\( x \), sedangkan yang kedua 
merupakan proyeksi pada sumbu-$y$.
\end{example}

% So not only is any linear map described by a
% matrix but any matrix describes a linear map.
Hasil ini berarti bahwa, jika praktis, kita dapat 
bekerja hanya dengan matriks,
melakukan perhitungan tanpa harus
mengkhawatirkan apakah matriks yang sedang dikaji merepresentasikan 
suatu peta linear pada sepasang ruang.

Ketika kita bekerja dengan suatu matriks 
tanpa membayangkan ruang atau basis tertentu,
kita dapat mengambil
domain dan kodomain sebagai $\Re^n$ dan $\Re^m$, dengan basis-basis
standar.
Ini praktis karena dengan basis standar,
representasi vektor bersifat transparan\Dash
representasi $\vec{v}$ adalah $\vec{v}$ sendiri.
(Dalam hal ini,
ruang kolom matriks sama dengan jangkauan peta dan,
akibatnya,
ruang kolom \( H \) sering dilambangkan dengan \( \rangespace{H} \).) 

% With the theorem, we have characterized that for fixed bases we can 
% pass back and forth between a linear map and its representation as a matrix.
% Each linear map is described by a matrix and each matrix describes a 
% linear map
Jika diberikan suatu matriks, untuk membentuk peta yang terkait kita dapat memilih dari banyak 
ruang domain dan kodomain, serta banyak basis untuk ruang-ruang itu. 
Jadi, satu matriks dapat merepresentasikan banyak peta.
Kita mengakhiri bagian ini dengan mengilustrasikan bagaimana matriks dapat 
memberikan informasi
tentang peta-peta yang terkait.

\begin{theorem} \label{th:RankMatEqRankMap}
%<*th:RankMatEqRankMap>
\index{rank}\index{homomorfisme!rank}\index{matriks!rank}
Rank suatu matriks sama dengan rank setiap peta yang direpresentasikannya.
%</th:RankMatEqRankMap>
\end{theorem}

\begin{proof}
%<*pf:RankMatEqRankMap0>
Andaikan matriks \( H \) berukuran \( \nbym{m}{n} \). 
Tetapkan ruang domain dan kodomain $V$ serta $W$ berdimensi $n$ dan~$m$ dengan
basis \( B=\sequence{\vec{\beta}_1,\dots,\vec{\beta}_n} \) dan \( D \).
Maka \( H \) merepresentasikan suatu peta linear $h$ antara ruang-ruang tersebut terhadap 
basis-basis ini, dengan ruang jangkauan
\begin{align*}
  \set{h(\vec{v})\suchthat \vec{v}\in V}
  &=\set{h(c_1\vec{\beta}_1+\dots+c_n\vec{\beta}_n)  
           \suchthat c_1,\dots,c_n\in\Re}                    \\
  &=\set{c_1h(\vec{\beta}_1)+\dots+c_nh(\vec{\beta}_n)
            \suchthat c_1,\dots,c_n\in\Re}
\end{align*}
adalah rentang $\spanof{\set{h(\vec{\beta}_1),\dots,h(\vec{\beta}_n)}}$.
Rank peta $h$ adalah dimensi ruang jangkauan ini.
%</pf:RankMatEqRankMap0>

%<*pf:RankMatEqRankMap1>
Rank matriks adalah dimensi ruang kolomnya,
yaitu rentang himpunan kolom-kolomnya
$\spanof{\,\set{\rep{h(\vec{\beta}_1)}{D},\dots,\rep{h(\vec{\beta}_n)}{D}}\,}$.
%</pf:RankMatEqRankMap1>

%<*pf:RankMatEqRankMap2>
Untuk melihat bahwa kedua rentang memiliki dimensi yang sama, ingat 
dari bukti Lema~I.\ref{lem:EqDimImpIso} bahwa jika kita menetapkan suatu basis, maka
representasi terhadap basis tersebut memberikan isomorfisme 
$\map{\mbox{Rep}_D}{W}{\Re^m}$.
Di bawah isomorfisme ini, terdapat hubungan linear di antara anggota-anggota
ruang jangkauan jika dan hanya jika hubungan yang sama berlaku dalam
ruang kolom, misalnya,
$\zero=c_1\cdot h(\vec{\beta}_1)+\dots+c_n\cdot h(\vec{\beta}_n)$ jika dan hanya jika
$\zero=c_1\cdot\rep{h(\vec{\beta}_1)}{D}+\dots+c_n\cdot\rep{h(\vec{\beta}_n)}{D}$.
Karena itu, suatu himpunan bagian ruang jangkauan bebas linear jika dan hanya jika
himpunan bagian yang bersesuaian dari ruang kolom bebas linear.
Oleh karena itu, ukuran himpunan bagian bebas linear terbesar dari
ruang jangkauan sama dengan ukuran himpunan bagian bebas linear terbesar dari
ruang kolom, sehingga kedua ruang memiliki dimensi yang sama.
%</pf:RankMatEqRankMap2>
\end{proof}

Ini menyelesaikan ambiguitas semu dalam penggunaan kata yang sama,
`rank', baik untuk matriks maupun peta.

\begin{example}
Setiap peta yang direpresentasikan oleh
\begin{equation*}
    \begin{mat}[r]
      1  &2  &2  \\
      1  &2  &1  \\
      0  &0  &3  \\
      0  &0  &2
    \end{mat}
\end{equation*}
harus memiliki domain berdimensi tiga dan kodomain berdimensi empat.
Selain itu, karena rank matriks ini adalah dua 
(kita dapat melihatnya langsung atau memperolehnya dengan Metode Gauss),
setiap peta yang direpresentasikan oleh matriks ini memiliki ruang jangkauan berdimensi dua.
\end{example}

\begin{corollary} \label{cor:MatDescsMap}
%<*co:MatDescsMap>
Misalkan $h$ adalah peta linear yang direpresentasikan oleh matriks $H$.
Maka $h$ bersifat pada jika dan hanya jika rank $H$ sama dengan banyaknya 
barisnya, 
dan $h$ bersifat satu-satu jika dan hanya jika rank $H$ sama dengan banyaknya
kolomnya.
%</co:MatDescsMap>
\end{corollary}

\begin{proof}
%<*pf:MatDescsMap0>
Untuk bagian sifat pada, dimensi ruang jangkauan $h$ adalah rank $h$,
yang menurut teorema sama dengan rank $H$.
Karena dimensi kodomain $h$ sama dengan banyaknya baris dalam $H$,
jika rank $H$ sama dengan banyaknya baris, maka dimensi
ruang jangkauan sama dengan dimensi kodomain.
Namun, subruang yang berdimensi sama dengan ruang induknya harus sama dengan
ruang induk tersebut
(karena sebarang basis ruang jangkauan merupakan himpunan bagian bebas linear 
dari kodomain
yang ukurannya sama dengan dimensi kodomain, sehingga 
basis ruang jangkauan ini juga harus merupakan 
basis kodomain).
%</pf:MatDescsMap0>

%<*pf:MatDescsMap1>
Untuk bagian lainnya, 
suatu peta linear bersifat satu-satu jika dan hanya jika peta tersebut merupakan isomorfisme
antara domain dan jangkauannya, yakni jika dan hanya jika domainnya memiliki
dimensi yang sama dengan jangkauannya.
Banyaknya kolom dalam $H$ adalah dimensi domain $h$ dan,
menurut teorema, rank $H$ sama dengan dimensi 
jangkauan $h$.
%</pf:MatDescsMap1>
\end{proof}

\begin{definition}  \label{df:NonsingularMap}
%<*df:NonsingularMap>
Pemetaan linear yang injektif dan surjektif disebut
\definend{nonsingular}\index{homomorfisme!nonsingular}%
\index{nonsingular!homomorfisme}; jika tidak, pemetaan itu disebut
\definend{singular}\index{homomorfisme!singular}%
\index{singular!homomorfisme}.
Dengan kata lain, pemetaan linear nonsingular jika dan hanya jika merupakan
isomorfisme.
%</df:NonsingularMap>
\end{definition}

\begin{remark}
Sebagian penulis memakai `nonsingular' sebagai sinonim untuk injektif,
sedangkan penulis lain memakainya seperti dalam buku ini.
Perbedaannya kecil karena setiap pemetaan bersifat surjektif pada daerah
hasilnya, sehingga pemetaan injektif merupakan isomorfisme dengan daerah hasilnya.
\end{remark}

Dalam bab pertama, kita mendefinisikan matriks nonsingular sebagai matriks
persegi yang menjadi matriks koefisien suatu sistem linear dengan solusi tunggal.
Hasil berikut membenarkan penggunaan ganda istilah tersebut.

\begin{lemma} \label{le:NonsingMatIffNonsingMap}
\index{homomorfisme!nonsingular}\index{matriks!nonsingular}
\index{nonsingular}
%<*le:NonsingMatIffNonsingMap>
Pemetaan linear nonsingular direpresentasikan oleh matriks persegi.
Matriks persegi merepresentasikan pemetaan nonsingular jika dan hanya jika
matriks itu nonsingular.
Jadi, suatu matriks merepresentasikan isomorfisme jika dan hanya jika matriks
itu persegi dan nonsingular.
%</le:NonsingMatIffNonsingMap>
\end{lemma}

\begin{proof}
%<*pf:NonsingMatIffNonsingMap0>
Andaikan pemetaan $\map{h}{V}{W}$ nonsingular.
\nearbycorollary{cor:MatDescsMap} menyatakan bahwa untuk setiap matriks $H$
yang merepresentasikan pemetaan tersebut, karena $h$ surjektif, banyak baris
$H$ sama dengan peringkat $H$; dan karena $h$ injektif, banyak kolom~$H$
juga sama dengan peringkat~$H$.
Jadi, $H$ persegi.
%</pf:NonsingMatIffNonsingMap0>

%<*pf:NonsingMatIffNonsingMap1>
Selanjutnya, andaikan $H$ persegi berukuran $\nbyn{n}$.
Matriks~$H$ nonsingular jika dan hanya jika peringkat barisnya~$n$;
menurut Teorema~Dua.III.\ref{th:RowRankEqualsColumnRank}, ini berlaku jika dan
hanya jika peringkat $H$ adalah~$n$; menurut
\nearbytheorem{th:RankMatEqRankMap}, ini berlaku jika dan hanya jika peringkat
$h$ adalah~$n$; dan menurut Teorema~I.\ref{th:NDimSpaceIsoRN}, ini berlaku jika
dan hanya jika $h$ merupakan isomorfisme.
(Pernyataan terakhir berlaku karena domain $h$ berdimensi~$n$, yaitu banyak
kolom dalam $H$.)
%</pf:NonsingMatIffNonsingMap1>
\end{proof}


\begin{example}
Setiap pemetaan dari \( \Re^2 \) ke \( \polyspace_1 \) yang terhadap
sebarang pasangan basis direpresentasikan oleh
\begin{equation*}
  \begin{mat}[r]
     1  &2  \\
     0  &3  
  \end{mat}
\end{equation*}
bersifat nonsingular karena matriks ini berperingkat dua.
\end{example}

\begin{example} \label{ex:NonSMatHasNonSMap}
Setiap pemetaan \( \map{g}{V}{W} \) yang direpresentasikan oleh
\begin{equation*}
  \begin{mat}[r]
    1  &2  \\
    3  &6
  \end{mat}
\end{equation*}
bersifat singular karena matriks ini singular.
\end{example}

Sekarang kita telah melihat bahwa hubungan antara pemetaan dan matriks berlaku
dua arah:~untuk suatu pasangan basis tertentu, setiap pemetaan linear
direpresentasikan oleh matriks dan setiap matriks mendeskripsikan suatu
pemetaan linear.
Dengan menetapkan ruang dan basis, kita memperoleh korespondensi antara
pemetaan dan matriks.
Pada sisa bab ini kita akan menjelajahi korespondensi tersebut.
Sebagai contoh, kita telah mendefinisikan operasi penjumlahan dan perkalian
skalar pada pemetaan linear, dan akan melihat operasi matriks yang bersesuaian.
Kita juga akan melihat operasi matriks yang merepresentasikan komposisi
pemetaan, serta cara mencari matriks yang merepresentasikan invers pemetaan.


\begin{exercises}
  \item Untuk setiap matriks, nyatakan dimensi domain dan kodomain setiap
    pemetaan yang direpresentasikannya.
    \begin{exparts*}
      \partsitem 
        $\begin{mat}
          2  &1  \\
          3  &4
        \end{mat}$
      \partsitem 
        $\begin{mat}
          1 &1 &-3 \\ 
          2 &5 &0
        \end{mat}$
      \partsitem 
        $\begin{mat}
           1 &3 \\
           1 &4 \\
           1 &-1
        \end{mat}$
      \partsitem 
        $\begin{mat}
           0 &0 &0 \\
           0 &0 &0
        \end{mat}$
      \partsitem 
        $\begin{mat}
           1 &-1 &4 &5 \\
           0 &0  &0 &0
        \end{mat}$
    \end{exparts*}
    \begin{answer}
      \begin{exparts}
        \partsitem domain: $2$, kodomain: $2$
        \partsitem domain: $3$, kodomain: $2$
        \partsitem domain: $2$, kodomain: $3$
        \partsitem domain: $3$, kodomain: $2$
        \partsitem domain: $4$, kodomain: $2$
      \end{exparts}
    \end{answer}
  \item Tinjau pemetaan linear $\map{f}{V}{W}$ yang terhadap suatu basis
    $B,D$ direpresentasikan oleh matriks berikut.
    Tentukan apakah pemetaan itu nonsingular.
    \begin{exparts*}
      \partsitem $\begin{mat}
           2  &1 \\
           3  &4
         \end{mat}$
      \partsitem $\begin{mat}
           1 &1 \\  
          -3 &-3
         \end{mat}$
      \partsitem $\begin{mat}
          3 &0 &0 \\
          2 &1 &0 \\
          4 &4 &4
         \end{mat}$
      \partsitem $\begin{mat}
          2 &0 &-2 \\
          1 &1 &0 \\
          4 &1 &-4
         \end{mat}$
    \end{exparts*}
    \begin{answer}
      Untuk setiap bagian, kita hanya perlu menentukan apakah matriksnya
      nonsingular, mungkin dengan Metode Gauss.
      (Sebenarnya, semua bagian dapat ditentukan hanya dengan mengamati.)
      \begin{exparts}
        \partsitem Matriks ini nonsingular karena baris kedua bukan kelipatan
          baris pertama, sehingga pemetaannya nonsingular.
        \partsitem Matriksnya singular, sehingga pemetaannya singular.
        \partsitem Nonsingular.
        \partsitem Nonsingular.
      \end{exparts}
    \end{answer}
  \recommended \item
   Misalkan $h$ adalah pemetaan linear yang didefinisikan oleh matriks berikut
   pada domain~$\polyspace_1$ dan kodomain~$\Re^2$ terhadap basis yang diberikan.
   \begin{equation*}
     H=
     \begin{mat}
       2  &1  \\
       4  &2
     \end{mat}
     \quad
     B=\sequence{1+x,x},\,
     D=\sequence{\colvec{1 \\ 1},\colvec{1 \\ 0}}
   \end{equation*}
   Apakah citra vektor $\vec{v}=2x-1$ oleh $h$?
   \begin{answer}
     Terhadap $B$, representasi vektornya adalah
     \begin{equation*}
       \rep{2x-1}{B}=\colvec{-1 \\ 3}
     \end{equation*}
     Dengan perkalian matriks-vektor, kita dapat menghitung $\rep{h(\vec{v})}{D}$.
     \begin{equation*}
       \rep{h(2x-1)}{D}
       =
       \begin{mat}
         2  &1  \\
         4  &2
       \end{mat}
       \colvec{-1 \\ 3}_B
       =
       \colvec{1 \\ 2}_D
     \end{equation*}
     Dari representasi tersebut, kita dapat menghitung $h(\vec{v})$.
     \begin{equation*}
       h(2x-1)=1\cdot\colvec{1 \\ 1}+2\cdot\colvec{1 \\ 0}
       =\colvec{3 \\ 1}
     \end{equation*}
   \end{answer}
 \recommended \item 
   Tentukan apakah setiap vektor berada dalam daerah hasil pemetaan dari
   \( \Re^3 \) ke \( \Re^2 \) yang terhadap basis standar direpresentasikan
   oleh matriks berikut.
   \begin{exparts*}
     \partsitem \( \begin{mat}[r]
                1  &1  &3  \\
                0  &1  &4
              \end{mat}  \),~\( \colvec[r]{1 \\ 3} \)
     \partsitem \( \begin{mat}[r]
                2  &0  &3  \\
                4  &0  &6
              \end{mat}  \),~\( \colvec[r]{1 \\ 1} \)
   \end{exparts*}
   \begin{answer} 
     Seperti dijelaskan dalam subbagian, terhadap basis standar representasi
     bersifat langsung. Jadi, misalnya, matriks pertama mendeskripsikan
     pemetaan berikut.
     \begin{equation*}
       \colvec[r]{1 \\ 0 \\ 0}=\colvec[r]{1 \\ 0 \\ 0}_{\stdbasis_3}
          \!\mapsto\colvec[r]{1 \\ 0}_{\stdbasis_2}=\colvec[r]{1 \\ 0}
       \qquad
       %\colvec[r]{0 \\ 1 \\ 0}=\colvec[r]{0 \\ 1 \\ 0}_{\stdbasis_3}
       %   \!\mapsto\colvec[r]{1 \\ 1}_{\stdbasis_2}=\colvec[r]{1 \\ 1}
       \colvec[r]{0 \\ 1 \\ 0}
          \!\mapsto\colvec[r]{1 \\ 1}
       \qquad
       %\colvec[r]{0 \\ 0 \\ 1}=\colvec[r]{0 \\ 0 \\ 1}_{\stdbasis_3}
       %   \!\mapsto\colvec[r]{3 \\ 4}_{\stdbasis_2}=\colvec[r]{3 \\ 4}
       \colvec[r]{0 \\ 0 \\ 1}
          \!\mapsto\colvec[r]{3 \\ 4}
     \end{equation*}
     Jadi, untuk bagian pertama kita menanyakan apakah terdapat skalar
     sedemikian sehingga
     \begin{equation*}
       c_1\colvec[r]{1 \\ 0}+c_2\colvec[r]{1 \\ 1}
           +c_3\colvec[r]{3 \\ 4}=\colvec[r]{1 \\ 3}
     \end{equation*}
     yaitu apakah vektor tersebut berada dalam ruang kolom matriks.
     \begin{exparts}
       \partsitem Ya.
         Kita dapat memperoleh kesimpulan ini dengan menyusun sistem linear
         yang dihasilkan dan menerapkan Metode Gauss seperti biasa.
         Cara lain adalah mengamati persamaan kolom dan mengambil
         $c_3=3/4$, $c_1=-5/4$, serta $c_2=0$.
         Cara ketiga adalah memperhatikan bahwa peringkat matriks sama dengan
         dua, yang sama dengan dimensi kodomain. Jadi pemetaan itu surjektif:
         daerah hasilnya seluruh $\Re^2$, termasuk vektor yang diberikan.
       \partsitem Tidak; perhatikan bahwa semua kolom matriks memiliki
         komponen kedua dua kali komponen pertama, sedangkan vektor yang
         diberikan tidak.
         Sebagai alternatif, ruang kolom matriks adalah
         \begin{equation*}
           \set{c_1\colvec[r]{2 \\ 4}
                +c_2\colvec[r]{0 \\ 0}
                +c_3\colvec[r]{3 \\ 6} \suchthat c_1,c_2,c_3\in\Re}
           =\set{c\colvec[r]{1 \\ 2}\suchthat c\in\Re}
         \end{equation*}
         (ini adalah fakta yang telah diamati, tetapi sekarang diperoleh lewat
         perhitungan, bukan ilham), dan vektor yang diberikan tidak berada
         dalam himpunan tersebut.
     \end{exparts}  
    \end{answer}
  \recommended \item  
    Tinjau matriks berikut, yang merepresentasikan transformasi $\Re^2$
    terhadap basis yang diberikan.
    \begin{equation*}
      \frac{1}{2}\cdot
      \begin{mat}[r]
        1  &1  \\
        -1 &1
      \end{mat}
      \qquad
      B=\sequence{\colvec[r]{0 \\ 1},\colvec[r]{1 \\ 0}}
      \quad 
      D=\sequence{\colvec[r]{1 \\ 1},\colvec[r]{1 \\ -1}}
    \end{equation*}
    \begin{exparts}
      \partsitem Ke vektor manakah dalam kodomain anggota pertama $B$ dipetakan?
      \partsitem Bagaimana dengan anggota kedua?
      \partsitem Ke manakah vektor umum dari domain, yaitu vektor dengan
        komponen $x$ dan $y$, dipetakan?
        Dengan kata lain, transformasi \( \Re^2 \) manakah yang
        direpresentasikan matriks ini terhadap \( B,D \)?
    \end{exparts}
    \begin{answer}
      \begin{exparts}
        \partsitem Anggota pertama basis
          \begin{equation*}
            \colvec[r]{0 \\ 1}=\colvec[r]{1 \\ 0}_B
          \end{equation*}
          dipetakan ke
          \begin{equation*}
            \colvec[r]{1/2 \\ -1/2}_D
          \end{equation*}
          yang merupakan anggota kodomain berikut.
          \begin{equation*}
            \frac{1}{2}\cdot\colvec[r]{1 \\ 1}
              -\frac{1}{2}\cdot\colvec[r]{1 \\ -1}
              =\colvec[r]{0 \\ 1}
          \end{equation*}
        \partsitem Anggota kedua basis dipetakan sebagai
          \begin{equation*}
            \colvec[r]{1 \\ 0}=\colvec[r]{0 \\ 1}_B
            \mapsto
            \colvec[r]{1/2 \\ 1/2}_D
          \end{equation*}
          ke anggota kodomain berikut.
          \begin{equation*}
            \frac{1}{2}\cdot\colvec[r]{1 \\ 1}
              +\frac{1}{2}\cdot\colvec[r]{1 \\ -1}
              =\colvec[r]{1 \\ 0}
          \end{equation*}
        \partsitem Karena pemetaan yang direpresentasikan matriks merupakan
          identitas pada basis, pemetaan itu harus menjadi identitas pada
          semua anggota domain.
          Kita dapat memperoleh kesimpulan yang sama dengan meninjau
          \begin{equation*}
            \colvec{x \\ y}=\colvec{y \\ x}_B
          \end{equation*}
          yang dipetakan ke
          \begin{equation*}
            \colvec{(x+y)/2 \\ (x-y)/2}_D
          \end{equation*}
          yang merepresentasikan anggota $\Re^2$ berikut.
          \begin{equation*}
            \frac{x+y}{2}\cdot\colvec[r]{1 \\ 1}
              +\frac{x-y}{2}\cdot\colvec[r]{1 \\ -1}
            =\colvec{x \\ y}
          \end{equation*}
      \end{exparts}
    \end{answer}
  \item Tinjau homomorfisme $\map{h}{\Re^2}{\Re^2}$ yang terhadap basis
    standar $\stdbasis_2,\stdbasis_2$ direpresentasikan oleh matriks berikut.
    \begin{equation*}
      \begin{mat}
        1  &3  \\
        2  &4
      \end{mat}
    \end{equation*}
    Carilah citra setiap vektor oleh $h$.
    \begin{exparts*}
      \partsitem $\colvec{2  \\ 3}$
      \partsitem $\colvec{0  \\ 1}$
      \partsitem $\colvec{-1 \\ 1}$
    \end{exparts*}
    \begin{answer}
      Karena terhadap basis standar~$\stdbasis_2$ suatu vektor
      direpresentasikan oleh dirinya sendiri, kita hanya perlu melakukan
      perkalian matriks-vektor.
      \begin{exparts}
        \partsitem 
          \begin{align*}
            \rep{h(\colvec{2 \\ 3})}{\stdbasis_2}       
             &=\rep{h}{\stdbasis_2,\stdbasis_2}\;
             \rep{\colvec{2 \\ 3}}{\stdbasis_2}        \\
             &=\begin{mat}
              1  &3  \\
              2  &4
             \end{mat}_{\stdbasis_2,\stdbasis_2}
             \colvec{2 \\ 3}_{\stdbasis_2}                 \\
             &=\colvec{11 \\ 16}_{\stdbasis_2}=
             \colvec{11 \\ 16}
          \end{align*}
        \partsitem
           $\begin{mat}
              1  &3  \\
              2  &4
             \end{mat}
             \colvec{0 \\ 1}=
             \colvec{3 \\ 4}$
        \partsitem
           $\begin{mat}
              1  &3  \\
              2  &4
             \end{mat}
             \colvec{-1 \\ 1}=
             \colvec{2 \\ 2}$
      \end{exparts}
    \end{answer}
  \item 
    Transformasi apakah pada
    \( F=\set{a\cos\theta+b\sin\theta\suchthat a,b\in\Re} \)
    yang terhadap \( B=\sequence{\cos\theta-\sin\theta,\sin\theta} \) dan
    \( D=\sequence{\cos\theta+\sin\theta,\cos\theta} \) direpresentasikan
    oleh matriks berikut?
    \begin{equation*}
      \begin{mat}[r]
          0  &0  \\
          1  &0
      \end{mat}
    \end{equation*}
    \begin{answer}
      Anggota umum domain direpresentasikan terhadap basis domain sebagai
      \begin{equation*}
        a\cos\theta+b\sin\theta=\colvec{a \\ a+b}_B
      \end{equation*}
      dan dipetakan ke
      \begin{equation*}
        \colvec{0 \\ a}_D
        \quad\text{yang merepresentasikan}\quad
        0\cdot(\cos\theta+\sin\theta)+a\cdot(\cos\theta)
      \end{equation*}
      sehingga pemetaan linear yang direpresentasikan matriks terhadap
      basis-basis tersebut,
      \begin{equation*}
        a\cos\theta+b\sin\theta  
            \mapsto
        a\cos\theta
      \end{equation*}
      merupakan proyeksi pada komponen pertama.
    \end{answer}
  \recommended \item  
     Tentukan apakah $1+2x$ berada dalam daerah hasil pemetaan dari $\Re^3$
     ke $\polyspace_2$ yang terhadap $\stdbasis_3$ dan
     $\sequence{1,1+x^2,x}$ direpresentasikan oleh matriks berikut.
     \begin{equation*}
       \begin{mat}[r]
         1  &3  &0  \\
         0  &1  &0  \\
         1  &0  &1
       \end{mat}
     \end{equation*}
     \begin{answer}
       Nyatakan basis $\polyspace_2$ yang diberikan dengan $B$.
       Penerapan pemetaan linear direpresentasikan oleh perkalian
       matriks-vektor.
       Jadi, vektor pertama dalam $\stdbasis_3$ dipetakan ke unsur
       $\polyspace_2$ yang terhadap~$B$ direpresentasikan oleh
       \begin{equation*}
       \begin{mat}[r]
         1  &3  &0  \\
         0  &1  &0  \\
         1  &0  &1
       \end{mat}
       \colvec[r]{1 \\ 0 \\ 0}
       =
       \colvec[r]{1 \\ 0 \\ 1}  
       \end{equation*}
       dan unsur tersebut adalah $1+x$.
       Hitung dua citra vektor basis lainnya dengan cara yang sama.
       \begin{equation*}
         \begin{mat}[r]
           1  &3  &0  \\
           0  &1  &0  \\
           1  &0  &1
         \end{mat}
         \colvec[r]{0 \\ 1 \\ 0}
         =
         \colvec[r]{3 \\ 1 \\ 0}=\rep{4+x^2}{B}
         \quad
         \begin{mat}[r]
           1  &3  &0  \\
           0  &1  &0  \\
           1  &0  &1
         \end{mat}
         \colvec[r]{0 \\ 0 \\ 1}
         =
         \colvec[r]{0 \\ 0 \\ 1}=\rep{x}{B}
       \end{equation*}
       Jadi, daerah hasil $h$ adalah rentang tiga polinomial
       $1+x$, $4+x^2$, dan $x$.
       Dengan demikian, kita dapat menentukan apakah $1+2x$ berada dalam
       daerah hasil pemetaan dengan mencari skalar $c_1$, $c_2$, dan $c_3$
       sedemikian sehingga
       \begin{equation*}
         c_1\cdot(1+x)+c_2\cdot(4+x^2)+c_3\cdot(x)=1+2x
       \end{equation*}
       dan jelas $c_1=1$, $c_2=0$, serta $c_3=1$ sudah cukup.
       Jadi, $1+2x$ berada dalam daerah hasil karena merupakan citra vektor
       berikut.
       \begin{equation*}
         1\cdot\colvec[r]{1 \\ 0 \\ 0}+0\cdot\colvec[r]{0 \\ 1 \\ 0}
             +1\cdot\colvec[r]{0 \\ 0 \\ 1}
       \end{equation*}

       \textit{Catatan.}
       Argumen yang lebih ringkas adalah memperhatikan bahwa matriks tersebut
       nonsingular, sehingga berperingkat~$3$ dan daerah hasilnya
       berdimensi~$3$. Karena kodomain juga berdimensi~$3$, pemetaan itu
       surjektif. Jadi setiap polinomial merupakan citra suatu vektor;
       khususnya,~$1+2x$ merupakan citra suatu vektor dalam domain.
     \end{answer}
  \item Carilah pemetaan yang direpresentasikan matriks berikut terhadap $B,B$.
    \begin{equation*}
      \begin{mat}
        2  &1  \\
        -1 &0
      \end{mat}
      \qquad
      B=\sequence{\colvec{1 \\ 0},\colvec{1 \\ 1}}
    \end{equation*}
    \begin{answer}
        Dengan $B=\sequence{\vec{\beta}_1,\vec{\beta}_2}$, kita dapat
          menentukan $\rep{\vec{v}}{B}$ melalui pengamatan, dengan $\vec{v}$
          sebagai vektor umum berentri $x$ dan~$y$.
          \begin{equation*}
            \colvec{x \\ y}=a\cdot\colvec{1 \\ 0}+b\cdot\colvec{1 \\ 1}
            \quad\text{menghasilkan}\quad
            b=y, a=x-y
          \end{equation*}
          Jadi, representasi vektor umum terhadap~$B$ adalah
          \begin{equation*}
            \rep{\colvec{x \\ y}}{B}=\colvec{x-y \\ y}
          \end{equation*}
          Hitung pengaruh pemetaan dengan perkalian matriks-vektor.
          \begin{equation*}
            \begin{mat}
              2  &1  \\
             -1 &0
            \end{mat}_B
            \colvec{x-y \\ y}_B
            =
            \colvec{2(x-y)+y \\ -(x-y)}_B     
            =
            \colvec{2x-y \\ -x+y}_B     
          \end{equation*}
          Selesaikan dengan mengubahnya kembali ke representasi vektor standar.
          \begin{equation*}
            \colvec{2x-y \\ -x+y}_B     
            =
            (2x-y)\cdot\colvec{1 \\ 0}+(-x+y)\cdot\colvec{1 \\ 1}
            =\colvec{x \\ -x+y}      
          \end{equation*}
    \end{answer}
  \item 
    \nearbyexample{ex:NonSMatHasNonSMap} memberikan matriks singular yang
    karena itu terkait dengan pemetaan-pemetaan singular.
    Kita tidak dapat menyatakan tindakan pemetaan terkait~$g$ pada unsur
    domain~$\vec{v}\in V$ karena kita tidak mengetahui domain~$V$,
    kodomain~$W$, ataupun basis awal dan akhir $B$ dan~$D$.
    Namun, kita dapat menghitung apa yang terjadi pada representasi
    $\rep{\vec{v}}{B}$.
    \begin{exparts}
      \partsitem Carilah himpunan vektor kolom yang merepresentasikan anggota
        ruang nol dari setiap pemetaan~$g$ yang direpresentasikan matriks ini.
      \partsitem Carilah nulitas setiap pemetaan~$g$ tersebut.
      \partsitem Carilah himpunan vektor kolom yang merepresentasikan anggota
        daerah hasil setiap pemetaan~$g$ yang direpresentasikan matriks ini.
      \partsitem Carilah peringkat setiap pemetaan~$g$ tersebut.
      \partsitem Periksa bahwa peringkat ditambah nulitas sama dengan dimensi domain.
    \end{exparts}
    \begin{answer}
      Misalkan matriks tersebut adalah $G$ dan merepresentasikan
      $\map{g}{V}{W}$ terhadap basis $B$ dan $D$.
      Karena $G$ memiliki dua kolom, domain $V$ berdimensi dua.
      Karena $G$ memiliki dua baris, kodomain $W$ berdimensi dua.
      Tindakan $g$ pada representasi $\rep{\vec{v}}{B}$ dari anggota umum
      domain adalah
      \begin{equation*}
        \colvec{x \\ y}_B 
         \;\mapsto\; 
        \colvec{x+2y \\ 3x+6y}_D
      \end{equation*}
      \begin{exparts}
        \partsitem Apa pun basis kodomain~$D$, satu-satunya representasi
           vektor nol $\zero_W$ adalah
           \begin{equation*}
             \rep{\zero}{D}=\colvec[r]{0 \\ 0}_D
           \end{equation*}
           sehingga himpunan representasi anggota ruang nol adalah
           \begin{equation*}
             \set{\colvec{x \\ y}_B\suchthat \text{$x+2y=0$ dan $3x+6y=0$}}
             =\set{y\cdot\colvec[r]{-2 \\ 1}_B\suchthat y\in\Re}
           \end{equation*}
         \partsitem Nulitasnya adalah~$1$.
           Pemetaan representasi $\map{\mbox{Rep}_B}{V}{\Re^2}$ dan inversnya
           merupakan isomorfisme, sehingga mempertahankan dimensi subruang.
           Subruang $\Re^2$ pada bagian sebelumnya berdimensi satu.
           Karena itu, citra subruang tersebut di bawah invers pemetaan
           representasi\Dash yakni ruang nol $g$\Dash juga berdimensi satu.
         \partsitem Himpunan representasi anggota daerah hasil adalah
           \begin{align*}
             \set{\colvec{x+2y \\ 3x+6y}_D\suchthat x,y\in\Re}
             &=\set{x\cdot\colvec{1 \\ 3}_D+y\cdot\colvec{2 \\ 6}_D\suchthat x,y\in\Re}  \\
             &=\set{k\cdot\colvec[r]{1 \\ 3}_D\suchthat k\in\Re}
           \end{align*}
         \partsitem Tentu saja, \nearbytheorem{th:RankMatEqRankMap}
           menyatakan bahwa peringkat pemetaan sama dengan peringkat matriks,
           yaitu satu. Sebagai alternatif, argumen yang sama seperti untuk
           ruang nol menunjukkan bahwa dimensi daerah hasil adalah satu.
         \partsitem Satu ditambah satu sama dengan dua.
      \end{exparts}
    \end{answer}
  \recommended \item
     Anggap setiap matriks merepresentasikan $\map{h}{\Re^m}{\Re^n}$
     terhadap basis standar.
     Untuk setiap matriks, (i)~nyatakan $m$ dan~$n$.
     Kemudian susun matriks diperbesar dengan matriks yang diberikan di kiri
     dan vektor yang merepresentasikan anggota daerah hasil di kanan
     (misalnya, jika kodomainnya~$\Re^3$, letakkan entri $a$, $b$, dan $c$
     pada kolom kanan).
     Lakukan reduksi Gauss--Jordan.
     Gunakan hasilnya untuk
     (ii)~mencari $\rangespace{h}$ dan $\rank(h)$ serta menyatakan apakah
     pemetaannya surjektif, dan
     (iii)~mencari $\nullspace{h}$ dan $\nullity(h)$ serta menyatakan apakah
     pemetaannya injektif.
     \begin{exparts}
       \partsitem     $
         \begin{mat}
           2  &1  \\ 
           -1 &3
         \end{mat}$
      \partsitem
        $\begin{mat}
          0  &1  &3  \\ 
          2  &3  &4  \\
         -2  &-1 &2 
        \end{mat}$
      \partsitem
        $\begin{mat}
          1  &1  \\ 
          2  &1  \\
          3  &1
       \end{mat}$
     \end{exparts}
     \begin{answer}
       \begin{exparts}
         \partsitem
           (i)~Dimensi ruang domain adalah banyaknya kolom,~$m=2$.
           Dimensi ruang kodomain adalah banyaknya baris,~$n=2$.

           Untuk bagian selanjutnya, tinjau persamaan matriks-vektor berikut.
           \begin{equation*}
             \begin{mat}
               2  &1  \\ 
              -1  &3        
             \end{mat}
             \colvec{x \\ y}
             =
             \colvec{a \\ b}
             \tag{$*$}
           \end{equation*}
           Kita menyelesaikannya untuk $x$ dan~$y$.
           \begin{equation*}
             \begin{amat}{2}
             2  &1 &a  \\ 
             -1 &3 &b       
             \end{amat}
             \grstep{(1/2)\rho_1+\rho_2}
             \grstep[(2/7)\rho_2]{(1/2)\rho_1}
             \grstep{-(1/2)\rho_2+\rho_1}
             \begin{amat}{2}
               1  &0 &(3/7)a-(1/7)b  \\ 
               0  &1 &(1/7)a+(2/7)b       
             \end{amat}
           \end{equation*}
           (ii)~Untuk setiap
           \begin{equation*}
             \colvec{a \\ b}\in\Re^2
           \end{equation*}
           dalam persamaan~($*$), perhitungan menunjukkan bahwa sistem
           memiliki solusi.
           Jadi, daerah hasil adalah seluruh kodomain,
           $\rangespace{h}=\Re^2$.
           Peringkat pemetaan adalah dimensi daerah hasil, yaitu~$2$.
           Pemetaan itu surjektif karena daerah hasilnya seluruh kodomain.

           (iii)~Sekali lagi berdasarkan perhitungan, untuk mencari ruang nol,
           tetapkan $a=b=0$ dalam persamaan~($*$). Hasilnya $x=y=0$.
           Ruang nol adalah subruang trivial dari domain.
           \begin{equation*}
             \nullspace{h}=\set{\colvec{0 \\ 0}}
           \end{equation*}
           Nulitas adalah dimensi ruang nol tersebut, yaitu~$0$.
           Pemetaan itu injektif karena ruang nolnya trivial.
         \partsitem
            (i)~Dimensi ruang domain adalah banyaknya kolom matriks,~$m=3$,
            sedangkan dimensi ruang kodomain adalah banyaknya baris,~$n=3$.

           Perhitungannya adalah sebagai berikut.
            \begin{multline*}
              \begin{amat}{3}
                0  &1  &3  &a \\ 
                2  &3  &4  &b \\
               -2  &-1 &2  &c
              \end{amat}
              \grstep{\rho_1\leftrightarrow\rho_2}
              \grstep{\rho_1+\rho_3}
              \grstep{-2\rho_2+\rho_3}                 \\
              \grstep{(1/2)\rho_1}
              \grstep{-(3/2)\rho_2+\rho_1}
              \begin{amat}{3}
                1  &0 &-5/2 &-(3/2)a+(1/2)b  \\ 
                0  &1 &3    &a  \\
                0  &0 &0    &-2a+b+c      
              \end{amat}
            \end{multline*}
            (ii)~Terdapat tripel kodomain
            \begin{equation*}
              \colvec{a \\ b \\ c}\in\Re^3
            \end{equation*}
            yang tidak menghasilkan solusi sistem; secara khusus, sistem
            hanya memiliki solusi jika $-2a+b+c=0$.
            \begin{equation*}
              \rangespace{h}=\set{\colvec{a \\ b \\ c}\suchthat a=(b+c)/2}
                            =\set{\colvec{1/2 \\ 1 \\ 0}b
                                  +\colvec{1/2 \\ 0 \\ 1}c\suchthat b,c\in\Re} 
            \end{equation*}
            Peringkat pemetaan adalah dimensi daerah hasil, yaitu~$2$.
            Pemetaan itu tidak surjektif karena daerah hasilnya bukan seluruh
            kodomain.

            (iii)~Menetapkan $a=b=c=0$ dalam perhitungan menghasilkan tak
            hingga banyak solusi.
            Parametrisasi dengan variabel bebas~$z$ memberikan deskripsi
            ruang nol berikut.
            \begin{equation*}
              \nullspace{h}=\set{\colvec{x \\ y \\ z}\suchthat 
                               \text{$y=-3z$ dan $x=(5/2)z$}}
                           =\set{\colvec{5/2 \\ -3 \\ 1}z\suchthat z\in\Re}
            \end{equation*}
            Nulitas adalah dimensi ruang nol tersebut, yaitu~$1$.
            Pemetaan itu tidak injektif karena ruang nolnya tidak trivial.
          \partsitem
            (i)~Domain berdimensi $m=2$, sedangkan kodomain berdimensi~$n=3$.
            Perhitungannya adalah
              \begin{equation*}
                \begin{amat}{2}
                  1  &1  &a \\ 
                  2  &1  &b \\
                  3  &1  &c
                \end{amat}
                \grstep[-3\rho_1+\rho_3]{-2\rho_1+\rho_2}
                \grstep{-2\rho_2+\rho_3}
                \grstep{-\rho_2}                 
                \grstep{-\rho_2+\rho_1}
                \begin{amat}{2}
                  1  &0 &-a+b  \\ 
                  0  &1 &2a-b  \\
                  0  &0 &a-2b+c      
                \end{amat}
              \end{equation*}

            (ii)~Daerah hasil adalah subruang kodomain berikut.
            \begin{equation*}
              \rangespace{h}=\set{\colvec{2b-c \\ b \\ c}\suchthat b,c\in\Re}
                  =\set{\colvec{2 \\ 1 \\ 0}b+\colvec{-1 \\ 0 \\ 1}c\suchthat b,c\in\Re}
            \end{equation*}
            Peringkatnya~$2$.
            Pemetaannya tidak surjektif.

            (iii)~Ruang nol adalah subruang trivial dari domain.
            \begin{equation*}
              \nullspace{h}=\set{\colvec{x \\ y}=\colvec{0 \\ 0}}
            \end{equation*}
            Nulitasnya~$0$.
            Pemetaannya injektif.
       \end{exparts}
     \end{answer}
  \item Gunakan metode dari latihan sebelumnya pada matriks berikut.
    \begin{equation*}
        \begin{mat}
          1 &0 &-1 \\
          2 &1 &0  \\
          2 &2 &2
        \end{mat}
    \end{equation*}
    \begin{answer}
          Berikut reduksi Gauss--Jordan-nya.
          \begin{align*}
            \begin{amat}{3}
                1 &0 &-1 &a \\
                2 &1 &0  &b \\
                2 &2 &2  &c  
            \end{amat}
            &\grstep[-2\rho_1+\rho_3]{-2\rho_1+\rho_2}
            \begin{amat}{3}
                1 &0 &-1 &a \\
                0 &1 &2  &-2a+b \\
                0 &2 &4  &-2a+c  
            \end{amat}                                    \\
            &\grstep{-2\rho_2+\rho_3}
            \begin{amat}{3}
                1 &0 &-1 &a \\
                0 &1 &2  &-2a+b \\
                0 &0 &0  &2a-2b+c  
            \end{amat}
          \end{align*}
          (i)~Dimensinya adalah $m=n=3$.
          (ii)~Daerah hasil adalah himpunan semua anggota kodomain yang
          membuat sistem ini memiliki solusi.
          \begin{equation*}
            \rangespace{h}=\set{\colvec{b-(1/2)c \\ b \\ c}\suchthat b,c\in\Re}
          \end{equation*}
          Peringkatnya adalah 2.
          Karena peringkat lebih kecil daripada dimensi kodomain~$n=3$,
          pemetaan itu tidak surjektif.
      
          (iii)~Ruang nol adalah himpunan anggota domain yang dipetakan ke
          $a=0$, $b=0$, dan~$c=0$.
          \begin{equation*}
            \nullspace{h}=\set{\colvec{z \\ -2z \\ z}\suchthat z\in\Re}
          \end{equation*}
          Nulitasnya~$1$.
          Karena nulitasnya tidak~$0$, pemetaan itu tidak injektif.
    \end{answer}
  \item Verifikasi bahwa pemetaan yang direpresentasikan matriks berikut
        merupakan isomorfisme.
        \begin{equation*}
          \begin{mat}
            2  &1  &0  \\
            3  &1  &1  \\
            7  &2  &1
          \end{mat}
        \end{equation*}
    \begin{answer}
       Untuk setiap pemetaan yang direpresentasikan matriks ini, domain dan
       kodomain masing-masing berdimensi~$3$.
           Untuk menunjukkan bahwa pemetaan tersebut merupakan isomorfisme,
           kita harus menunjukkan bahwa ia surjektif sekaligus injektif.
           Untuk itu, kita tidak perlu memperbesar matriks dengan $a$, $b$,
           dan~$c$; perhitungan berikut
             \begin{multline*}
               \begin{mat}
                 2  &1  &0 \\ 
                 3  &1  &1 \\
                 7  &2  &1
               \end{mat}
               \grstep[-(7/2)\rho_1+\rho_3]{-(3/2)\rho_1+\rho_2}
               \grstep{-3\rho_2+\rho_3}
               \grstep[-2\rho_2 \\ -(1/2)\rho_3]{(1/2)\rho_1}                 
               \grstep{2\rho_3+\rho_2}                          \\
               \grstep{-(1/2)\rho_2+\rho_1}
               \begin{mat}
                 1  &0 &0  \\ 
                 0  &1 &0  \\
                 0  &0 &1      
               \end{mat}
             \end{multline*}
          menunjukkan bahwa setiap vektor kodomain terkait dengan tepat satu
          vektor domain.
    \end{answer}
  \item Berikut adalah pembuktian alternatif untuk
    \nearbylemma{le:NonsingMatIffNonsingMap}.
    Diberikan matriks $\nbyn{n}$ $H$, tetapkan domain~$V$ dan kodomain~$W$
    dengan dimensi~$n$ yang sesuai, serta basis $B,D$ bagi ruang-ruang itu,
    lalu tinjau pemetaan~$h$ yang direpresentasikan matriks tersebut.
    \begin{exparts}
      \partsitem
        Tunjukkan bahwa $h$ surjektif jika dan hanya jika terdapat sedikitnya
        satu $\rep{\vec{v}}{B}$ yang oleh $H$ dikaitkan dengan setiap
        $\rep{\vec{w}}{D}$.
      \partsitem
        Tunjukkan bahwa $h$ injektif jika dan hanya jika terdapat paling banyak
        satu $\rep{\vec{v}}{B}$ yang oleh $H$ dikaitkan dengan setiap
        $\rep{\vec{w}}{D}$.
      \partsitem
        Tinjau sistem linear
        $H\cdot\rep{\vec{v}}{B}=\rep{\vec{w}}{D}$.
        Tunjukkan bahwa $H$ nonsingular jika dan hanya jika terdapat tepat satu
        solusi~$\rep{\vec{v}}{B}$ untuk setiap $\rep{\vec{w}}{D}$.
    \end{exparts}
    \begin{answer}
      \begin{exparts}
        \partsitem
          Pemetaan $h$ yang didefinisikan bersifat surjektif jika dan hanya
          jika untuk setiap $\vec{w}\in W$ terdapat $\vec{v}\in V$ sedemikian
          sehingga $h(\vec{v})=\vec{w}$.
          Karena setiap vektor memiliki tepat satu representasi, pengubahan ke
          representasi menunjukkan bahwa $h$ surjektif jika dan hanya jika
          untuk setiap representasi $\rep{\vec{w}}{D}$ terdapat representasi
          $\rep{\vec{v}}{B}$ sedemikian sehingga
          $H\cdot\rep{\vec{v}}{B}=\rep{\vec{w}}{D}$.
       \partsitem
          Bagian ini sama seperti bagian sebelumnya.
       \partsitem 
          Seperti dijelaskan pada awal subbagian ini, menurut definisi
          pemetaan $h$ yang ditentukan oleh matriks $H$ mengaitkan vektor
          domain $\vec{v}$ berikut dengan vektor kodomain $\vec{w}$ berikut.
          \begin{equation*}
            \rep{\vec{v}}{B}=\colvec{v_1 \\ \vdots \\ v_n}
            \qquad
            \rep{\vec{w}}{D}
               =
               H\cdot\rep{\vec{v}}{B}
               =\colvec{h_{1,1}v_1+\dots+h_{1,n}v_n  \\ 
                          \vdots \\ 
                       h_{m,1}v_1+\dots+h_{m,n}v_n}
          \end{equation*}
          Tetapkan $\vec{w}\in W$ dan tinjau sistem linear yang didefinisikan
          oleh persamaan di atas.
          \begin{equation*}
            % H\cdot\rep{\vec{v}}{B}=\rep{\vec{w}}{D}
            % \qquad
            \begin{linsys}{3}
              h_{1,1}v_1 &+ &\cdots &+ &h_{1,n}v_n &= &w_1  \\
              h_{2,1}v_1 &+ &\cdots &+ &h_{2,n}v_n &= &w_2  \\
                       &  &        & &          &\vdotswithin{=} & \\
              h_{n,1}v_1 &+ &\cdots &+ &h_{n,n}v_n &= &w_n
            \end{linsys}
          \end{equation*}
          (Sekali lagi, di sini $w_i$ ditetapkan dan $v_j$ adalah variabel tak
           diketahui.)
          Matriks $H$ nonsingular jika dan hanya jika untuk semua
          $w_1$, \ldots, $w_n$ sistem ini memiliki solusi tunggal.
          Berdasarkan dua bagian pertama latihan ini, hal tersebut berlaku
          jika dan hanya jika pemetaan $h$ surjektif dan injektif.
          Pada gilirannya, ini berlaku jika dan hanya jika $h$ merupakan
          isomorfisme.
       \end{exparts}
      \end{answer}
  \recommended \item  
    Karena peringkat matriks sama dengan peringkat setiap pemetaan yang
    direpresentasikannya, jika satu matriks merepresentasikan dua pemetaan berbeda,
    \( H=\rep{h}{B,D}=\rep{\hat{h}}{\hat{B},\hat{D}} \) 
    (dengan \( \map{h,\hat{h}}{V}{W} \))
    maka dimensi daerah hasil \( h \) sama dengan dimensi daerah hasil
    \( \hat{h} \).
    Apakah kedua daerah hasil yang berdimensi sama ini harus benar-benar sama?
    \begin{answer}
      Tidak, daerah hasilnya dapat berbeda.
      \nearbyexample{ex:CngBasesChgMap} menunjukkan hal ini.
    \end{answer}
  \item 
    Misalkan \( V \) adalah ruang berdimensi \( n \) dengan basis \( B \)
    dan \( D \).
    Tinjau pemetaan yang untuk \( \vec{v}\in V \) mengirim vektor kolom yang
    merepresentasikan \( \vec{v} \) terhadap \( B \) ke vektor kolom yang
    merepresentasikan \( \vec{v} \) terhadap \( D \).
    Tunjukkan bahwa pemetaan tersebut merupakan transformasi linear \( \Re^n \).
    \begin{answer}
      Ingat bahwa pemetaan representasi
      \begin{equation*}
        V\mapsunder{\text{Rep}_{B}}\Re^n
      \end{equation*}
      merupakan isomorfisme.
      Jadi, pemetaan inversnya $\map{\mbox{Rep}_B^{-1}}{\Re^n}{V}$ juga
      merupakan isomorfisme.
      Transformasi $\Re^n$ yang diinginkan kemudian merupakan komposisi berikut.
      \begin{equation*}
        \Re^n\mapsunder{\text{Rep}_{B}^{-1}}
        V\mapsunder{\text{Rep}_{D}}\Re^n
      \end{equation*}
      Karena komposisi isomorfisme juga merupakan isomorfisme, pemetaan
      $\composed{\mbox{Rep}_{D}}{\mbox{Rep}_{B}^{-1}}$ ini merupakan
      isomorfisme.
    \end{answer}
  \item 
    \nearbyexample{ex:CngBasesChgMap} menunjukkan bahwa mengubah pasangan
    basis dapat mengubah pemetaan yang direpresentasikan suatu matriks,
    walaupun domain dan kodomainnya tetap sama.
    Dapatkah pemetaannya tidak berubah?
    Apakah terdapat matriks \( H \), ruang vektor \( V \) dan \( W \), serta
    pasangan basis terkait \( B_1,D_1 \) dan \( B_2,D_2 \), dengan
    \( B_1\neq B_2 \), \( D_1\neq D_2 \), atau keduanya, sedemikian sehingga
    pemetaan yang direpresentasikan \( H \) terhadap \( B_1,D_1 \) sama
    dengan pemetaan yang direpresentasikannya terhadap \( B_2,D_2 \)?
    \begin{answer}
      Ya.
      Tinjau
      \begin{equation*}
        H=\begin{mat}[r]
            1  &0  \\
            0  &1
          \end{mat}
      \end{equation*}
      yang merepresentasikan pemetaan dari \( \Re^2 \) ke \( \Re^2 \).
      Terhadap basis standar \( B_1=\stdbasis_2, D_1=\stdbasis_2 \), matriks
      ini merepresentasikan pemetaan identitas.
      Terhadap
      \begin{equation*}
        B_2=D_2=\sequence{\colvec[r]{1 \\ 1},\colvec[r]{1 \\ -1}}
      \end{equation*}
      matriks ini kembali merepresentasikan identitas.
      Bahkan, selama basis awal dan akhir sama\Dash selama
      \( B_i=D_i \)\Dash pemetaan yang direpresentasikan $H$ adalah identitas.
    \end{answer}
  \recommended \item 
    Matriks persegi disebut matriks
    \definend{diagonal}\index{matriks!diagonal}\index{matriks diagonal} %
    jika semua entrinya nol kecuali mungkin entri pada diagonal dari kiri atas
    ke kanan bawah\Dash entri \( 1,1 \), entri \( 2,2 \), dan seterusnya.
    Tunjukkan bahwa pemetaan linear merupakan isomorfisme jika terdapat basis
    sedemikian sehingga terhadap basis tersebut pemetaan direpresentasikan
    oleh matriks diagonal tanpa entri nol pada diagonalnya.
    \begin{answer}
      Hal ini langsung mengikuti \nearbylemma{le:NonsingMatIffNonsingMap}.
    \end{answer}
  \item 
    Deskripsikan secara geometris tindakan pada \( \Re^2 \) dari pemetaan
    yang terhadap basis standar $\stdbasis_2,\stdbasis_2$ direpresentasikan
    oleh matriks berikut.
    \begin{equation*}
      \begin{mat}[r]
        3  &0  \\
        0  &2
      \end{mat}
    \end{equation*}
    Lakukan hal yang sama untuk matriks-matriks berikut.
    \begin{equation*}
      \begin{mat}[r]
        1  &0  \\
        0  &0
      \end{mat}
      \quad
      \begin{mat}[r]
        0  &1  \\
        1  &0
      \end{mat}
      \quad
      \begin{mat}[r]
        1  &3  \\
        0  &1
      \end{mat}
    \end{equation*}
    \begin{answer}
      Pemetaan pertama
      \begin{equation*}
        \colvec{x \\ y}=\colvec{x \\ y}_{\stdbasis_2}
        \mapsto
        \colvec{3x \\ 2y}_{\stdbasis_2}=\colvec{3x \\ 2y}
      \end{equation*}
      meregangkan vektor dengan faktor tiga pada arah~\( x \) dan faktor dua
      pada arah~\( y \).
      Pemetaan kedua
      \begin{equation*}
        \colvec{x \\ y}=\colvec{x \\ y}_{\stdbasis_2}
        \mapsto
        \colvec{x \\ 0}_{\stdbasis_2}=\colvec{x \\ 0}
      \end{equation*}
      memproyeksikan vektor pada sumbu~\( x \).
      Pemetaan ketiga
      \begin{equation*}
        \colvec{x \\ y}=\colvec{x \\ y}_{\stdbasis_2}
        \mapsto
        \colvec{y \\ x}_{\stdbasis_2}=\colvec{y \\ x}
      \end{equation*}
      menukar komponen pertama dan kedua (yakni refleksi terhadap garis
      \( y=x \)).
      Pemetaan terakhir
      \begin{equation*}
        \colvec{x \\ y}=\colvec{x \\ y}_{\stdbasis_2}
        \mapsto
        \colvec{x+3y \\ y}_{\stdbasis_2}=\colvec{x+3y \\ y}
      \end{equation*}
      menggeser vektor sejajar sumbu~\( x \) sejauh tiga kali jaraknya dari
      sumbu tersebut (ini adalah suatu \definend{geseran miring}).
     \end{answer}
  \item 
     Fakta bahwa untuk setiap pemetaan linear, peringkat ditambah nulitas sama
     dengan dimensi domain menunjukkan bahwa syarat perlu bagi keberadaan
     homomorfisme surjektif dari suatu ruang ke ruang kedua adalah tidak
     terjadinya kenaikan dimensi.
     Dengan kata lain, jika $\map{h}{V}{W}$ surjektif, dimensi $W$ harus lebih
     kecil daripada atau sama dengan dimensi $V$.
     \begin{exparts}
       \partsitem Tunjukkan bahwa kebalikan kuat berikut berlaku:
          tidak adanya kenaikan dimensi mengimplikasikan bahwa terdapat suatu
          homomorfisme; bahkan setiap matriks dengan ukuran dan peringkat yang
          tepat merepresentasikan pemetaan seperti itu.
       \partsitem Apakah terdapat basis untuk $\Re^3$ sedemikian sehingga
          matriks berikut
          \begin{equation*}
            H=\begin{mat}[r]
                1  &0  &0 \\
                2  &0  &0 \\
                0  &1  &0 
              \end{mat}
          \end{equation*}
          merepresentasikan pemetaan dari $\Re^3$ ke $\Re^3$ yang daerah
          hasilnya adalah subruang bidang~$xy$ dari $\Re^3$?
     \end{exparts}
     \begin{answer}
       \begin{exparts}
         \partsitem Hal ini langsung mengikuti
           \nearbytheorem{th:RankMatEqRankMap}.
         \partsitem Ya.
            Hal ini langsung mengikuti bagian sebelumnya.

            Untuk memberikan contoh khusus, kita dapat memulai dengan
            $\stdbasis_3$ sebagai basis domain, lalu mencari basis $D$ bagi
            kodomain $\Re^3$.
            Matriks $H$ memberikan tindakan pemetaan sebagai berikut.
            \begin{multline*}
              \colvec[r]{1 \\ 0 \\ 0}=\colvec[r]{1 \\ 0 \\ 0}_{\stdbasis_3}
                 \mapsto\colvec[r]{1 \\ 2 \\ 0}_D
              \quad       
              \colvec[r]{0 \\ 1 \\ 0}=\colvec[r]{0 \\ 1 \\ 0}_{\stdbasis_3}
                 \mapsto\colvec[r]{0 \\ 0 \\ 1}_D
              \\  %\quad       
              \colvec[r]{0 \\ 0 \\ 1}=\colvec[r]{0 \\ 0 \\ 1}_{\stdbasis_3}
                 \mapsto\colvec[r]{0 \\ 0 \\ 0}_D
            \end{multline*}
            Kita dapat mencari basis $D$ sedemikian sehingga
            \begin{equation*}
              \rep{\colvec[r]{1 \\ 0 \\ 0}}{D}=\colvec[r]{1 \\ 2 \\ 0}_D
              \quad\mbox{dan}\quad
              \rep{\colvec[r]{0 \\ 1 \\ 0}}{D}=\colvec[r]{0 \\ 0 \\ 1}_D
            \end{equation*}
            yakni sedemikian sehingga pemetaan yang direpresentasikan $H$
            terhadap $\stdbasis_3,D$ merupakan proyeksi pada bidang~$xy$.
            Syarat kedua memberikan anggota ketiga $D$ sebagai $\vec{e}_2$.
            Syarat pertama menyatakan bahwa anggota pertama $D$ ditambah dua
            kali anggota kedua sama dengan $\vec{e}_1$, sehingga basis berikut
            memenuhi.
            \begin{equation*}
              D=\sequence{\colvec{2 \\ 0 \\ 1},
                          \colvec{-1/2 \\ 0 \\ -1/2},
                          \colvec{0 \\ 1 \\ 0}}
            \end{equation*}
% sage: load('gauss_method.sage')
% sage: M = matrix(QQ, [[2,-1/2,0], [0,0,1], [1,-1/2,0]])
% sage: gauss_method(M)
% [   2 -1/2    0]
% [   0    0    1]
% [   1 -1/2    0]
% (' take', -1/2, 'times row', 1, 'plus row', 3)
% [   2 -1/2    0]
% [   0    0    1]
% [   0 -1/4    0]
% (' swap row', 2, 'with row', 3)
% [   2 -1/2    0]
% [   0 -1/4    0]
% [   0    0    1]
       \end{exparts} 
     \end{answer}
  \item 
    Misalkan \( V \) adalah ruang berdimensi \( n \) dan
    \( \vec{x}\in\Re^n \).
    Tetapkan basis \( B \) bagi \( V \), lalu tinjau pemetaan
    \( \map{h_{\vec{x}}}{V}{\Re} \) yang diberikan oleh hasil kali titik
    $\vec{v}\mapsto\vec{x}\dotprod\rep{\vec{v}}{B}$.
    \begin{exparts}
      \partsitem Tunjukkan bahwa pemetaan ini linear.
      \partsitem Tunjukkan bahwa untuk setiap pemetaan linear
        \( \map{g}{V}{\Re} \), terdapat \( \vec{x}\in\Re^n \) sedemikian
        sehingga \( g=h_{\vec{x}} \).
      \partsitem Pada bagian sebelumnya kita menetapkan basis dan memvariasikan
        \( \vec{x} \) untuk memperoleh semua pemetaan linear yang mungkin.
        Dapatkah kita memperoleh semua pemetaan linear dengan menetapkan
        \( \vec{x} \) dan memvariasikan basis?
    \end{exparts}
    \begin{answer}
      \begin{exparts}
        \partsitem Ingat bahwa pemetaan representasi
          $\map{\mbox{Rep}_{B}}{V}{\Re^n}$ bersifat linear (sebenarnya
          merupakan isomorfisme, tetapi sifat injektif atau surjektifnya tidak
          diperlukan di sini).
          Dengan menganggap transpos vektor kolom $\vec{x}$ sebagai matriks
          $\nbym{1}{n}$, pemetaan dari $\Re^n$ ke $\Re$ yang membawa vektor
          kolom ke hasil kali titiknya dengan $\vec{x}$ bersifat linear
          (ini adalah perkalian matriks-vektor, sehingga
          \nearbytheorem{th:MatIsLinMap} berlaku).
          Jadi, pemetaan $h_{\vec{x}}$ linear karena merupakan komposisi dua
          pemetaan linear.
          \begin{equation*}
            \vec{v}\mapsto \rep{\vec{v}}{B}
                   \mapsto \vec{x}\cdot\rep{\vec{v}}{B}     
          \end{equation*}
       \partsitem Setiap pemetaan linear $\map{g}{V}{\Re}$ direpresentasikan
          oleh suatu matriks
          \begin{equation*}
            \begin{mat}
              g_1  &g_2 &\cdots &g_n
            \end{mat}
          \end{equation*}
          (matriks tersebut memiliki $n$ kolom karena $V$ berdimensi~$n$ dan
          hanya satu baris karena $\Re$ berdimensi satu).
          Ambil $\vec{x}$ sebagai vektor kolom yang merupakan transpos matriks ini,
          \begin{equation*}
            \vec{x}=\colvec{g_1 \\ \vdots \\ g_n}
          \end{equation*}
          sehingga diperoleh tindakan yang diinginkan.
          \begin{equation*}
            \vec{v}=\colvec{v_1 \\ \vdots \\ v_n}
             \mapsto
            \colvec{g_1 \\ \vdots \\ g_n}\dotprod\colvec{v_1 \\ \vdots \\ v_n}
            =g_1v_1+\dots+g_nv_n
          \end{equation*}
        \partsitem Tidak.
          Jika \( \vec{x} \) memiliki suatu entri tak nol, maka
          \( h_{\vec{x}} \) tidak mungkin menjadi pemetaan nol; jika
          \( \vec{x} \) adalah vektor nol, maka \( h_{\vec{x}} \) hanya dapat
          menjadi pemetaan nol.
      \end{exparts}  
    \end{answer}
  \item
    Misalkan \( V,W,X \) adalah ruang vektor dengan basis \( B,C,D \).
    \begin{exparts}
      \partsitem Andaikan \( \map{h}{V}{W} \) direpresentasikan terhadap
        \( B,C \) oleh matriks \( H \).
        Berikan matriks yang merepresentasikan kelipatan skalar \( rh \),
        dengan \( r\in\Re \), terhadap \( B,C \), dan nyatakan dalam $H$.
      \partsitem Andaikan \( \map{h,g}{V}{W} \) direpresentasikan terhadap
        \( B,C \) oleh \( H \) dan \( G \).
        Berikan matriks yang merepresentasikan \( h+g \) terhadap \( B,C \),
        dan nyatakan dalam \( H \) serta \( G \).
      \partsitem Andaikan \( \map{h}{V}{W} \) direpresentasikan terhadap
        \( B,C \) oleh \( H \), dan \( \map{g}{W}{X} \) direpresentasikan
        terhadap \( C,D \) oleh \( G \).
        Berikan matriks yang merepresentasikan \( \composed{g}{h} \) terhadap
        \( B,D \), dan nyatakan dalam \( H \) serta \( G \).
    \end{exparts}
    \begin{answer}
       Lihat bagian berikutnya.
    \end{answer}
\end{exercises}
